\BourbakiExercisesSection

\begin{enumerate}
\def\labelenumi{\arabic{enumi}.}
\item
  Enumerate the elements of length \(\leqslant 4\) in \(\mathbf{M}(\mathbf{X})\) when \(\mathbf{X}\) consists of a single element.
\item
  Let X be a set and M one of the magmas M(X), Mo(X) or \(\mathrm{N}^{(X)}\). Show that every automorphism of M leaves X stable. Deduce an isomorphism of \(\mathfrak{S}_{X}\) onto the group Aut(M).
\end{enumerate}

Show that the endomorphisms of M correspond bijectively to the mappings of X into M.

\begin{enumerate}
\def\labelenumi{\arabic{enumi}.}
\setcounter{enumi}{2}
\tightlist
\item
  Let X be a set and \(M_{a}\) the free magma constructed on a set \(\{a\}\) with one element. Let \(\rho\) denote the homomorphism of \(\mathrm{M}(\mathrm{X})\) into \(M_{a}\) which maps every element of X onto a; on the other hand let \(\lambda: \mathrm{M}(\mathrm{X}) \to \mathrm{Mo}(\mathrm{X})\) be the homomorphism defined in no. 9. Show that the mapping \(w \mapsto (\lambda(w), \rho(w))\) is an isomorphism of \(\mathrm{M}(\mathrm{X})\) onto the submagma of \(\mathrm{Mo}(\mathrm{X}) \times \mathrm{M}_{a}\) consisting of the ordered pairs \((u, v)\) where u and v have the same length.
\end{enumerate}

¶ 4. Let X be a set and c an element not belonging to X. We write Y = X ∪ \{c\}. Let M be the magma with underlying set Mo(Y) and law of composition given by (u, v) ↦ cuv. Let L be the submagma of M generated by X and let ε be the mapping from Y to Z equal to -1 on X and to 1 at c.

\begin{enumerate}
\def\labelenumi{(\alph{enumi})}
\tightlist
\item
  Show that L is the subset of M consisting of the words \(w = y_{1} \ldots y_{p}\) , \(y_{i} \in Y\) , satisfying the two following conditions:
\end{enumerate}

\[
(i) \varepsilon (y _ {1}) + \dots + \varepsilon (y _ {p}) = - 1
\]

\[
(ii) \varepsilon \left(y _ {1}\right) + \dots + \varepsilon \left(y _ {i}\right) \geqslant 0 \quad \text { for } 1 \leqslant i \leqslant p - 1.
\]

Show that every element of L - X may be written uniquely in the form cuv with \(u, v \in L\) .

(The elements \(w\) satisfying (i) and (ii) form a submagma \(L'\) of \(M\) containing \(X\) . If \(w = y_1 \ldots y_p\) belongs to \(L'\) , let \(k\) be the least integer such that \(\varepsilon(y_1) + \cdots + \varepsilon(y_k) = 0\) ; show that \(y_1 = c\) and that the elements \(u = y_2 \ldots y_k\) and \(v = y_{k+1} \ldots y_p\) belong to \(L'\) ; then \(w = cuv\) , whence by induction on \(p\) the relation \(w \in L\) . Then show that the relations \(u', v' \in L\) and \(w = cu'v'\) imply \(u' = u, v' = v\) .)

\begin{enumerate}
\def\labelenumi{(\alph{enumi})}
\setcounter{enumi}{1}
\tightlist
\item
  Show that the injection of X into L extends to an isomorphism of M(X) onto L.
\end{enumerate}

\begin{enumerate}
\def\labelenumi{\arabic{enumi}.}
\setcounter{enumi}{4}
\tightlist
\item
  Let X be a set with one element. If n is an integer \(\geqslant1\) , let \(u_{n}\) denote the number of elements in M(X) of length n.
\end{enumerate}

\begin{enumerate}
\def\labelenumi{(\alph{enumi})}
\tightlist
\item
  For every set Y, show that
\end{enumerate}

\[
\operatorname{Card} \left(\mathrm{M} _ {n} (\mathrm{Y})\right) = u _ {n}. \operatorname{Card} (\mathrm{Y}) ^ {n}.
\]

(Use Exercise 3.)

\begin{enumerate}
\def\labelenumi{(\alph{enumi})}
\setcounter{enumi}{1}
\tightlist
\item
  Establish the relation
\end{enumerate}

\[
u _ {n} = \sum_ {p = 1} ^ {n - 1} u _ {p} u _ {n - p} \quad \text { for } n \geqslant 2.
\]

\begin{enumerate}
\def\labelenumi{(\alph{enumi})}
\setcounter{enumi}{2}
\tightlist
\item
  *Let \(f(\mathbf{T})\) be the formal power series \(\sum_{n=1}^{\infty} u_n \mathbf{T}^n\) . Show that
\end{enumerate}

\[
f (\mathbf {T}) = \mathbf {T} + f (\mathbf {T}) ^ {2}.
\]

Deduce the formulae

\[
f (\mathrm{T}) = \frac {1}{2} (1 - (1 - 4 \mathrm{T}) ^ {1 / 2}), \quad u _ {n} = \frac {2 ^ {n - 1}}{n !} 1. 3. 5 \dots (2 n - 3) \quad \text { for } n \geqslant 2. *
\]

Obtain the last result using Exercise 11 of Set Theory, III, § 5.

¶ 6. Let X be a set.

\begin{enumerate}
\def\labelenumi{(\alph{enumi})}
\item
  Let N be a submagma of M(X) and Y = N - (N.N). The injection Y → N extends to a homomorphism u: M(Y) → N. Show that u is an isomorphism. (In other words, every submagma of a free magma is free.)
\item
  If \(x \in \mathbf{M}(\mathbf{X})\) , let \(\mathbf{M}_x\) denote the submagma of \(\mathbf{M}(\mathbf{X})\) generated by \(x\) ; by (a), \(\mathbf{M}_x\) is identified with the free magma constructed on \(x\) . Show that, if \(x, y \in \mathbf{M}(\mathbf{X})\) , then either \(\mathbf{M}_x \subset \mathbf{M}_y\) or \(\mathbf{M}_y \subset \mathbf{M}_x\) or \(\mathbf{M}_x \cap \mathbf{M}_y = \varnothing\) . (If \(\mathbf{M}_x \cap \mathbf{M}_y \neq \varnothing\) , let \(z\) be an element of \(\mathbf{M}_x \cap \mathbf{M}_y\) of minimum length; show that either \(z = x\) or \(z = y\) .)
\end{enumerate}

\begin{enumerate}
\def\labelenumi{\arabic{enumi}.}
\setcounter{enumi}{6}
\tightlist
\item
  Let \(\mathbf{M}_x\) be the free magma constructed on a set \(\{x\}\) with one element. (a) For all \(y \in \mathbf{M}_x\) , show that there exists a unique endomorphism \(f_y\) of \(\mathbf{M}_x\) such that \(f_y(x) = y\) ; it is an isomorphism of \(\mathbf{M}_x\) onto the submagma \(\mathbf{M}_y\) generated by \(y\) (cf.~Exercise 6).
\end{enumerate}

\begin{enumerate}
\def\labelenumi{(\alph{enumi})}
\setcounter{enumi}{1}
\item
  If \(y, z \in \mathbf{M}_x\) , we write \(y \circ z = f_y(z)\) . Show that the law of composition \((y, z) \mapsto y \circ z\) makes \(\mathbf{M}_x\) into a monoid with identity element \(x\) which is isomorphic to the monoid \(\operatorname{End}(\mathbf{M}_x)\) of endomorphisms of \(\mathbf{M}_x\) .
\item
  Let \(\mathbf{N}\) be the set of monogenous submagmas of \(\mathbf{M}_x\) distinct from \(\mathbf{M}_x\) . An element \(y\) of \(\mathbf{M}_x\) is called primitive if \(\mathbf{M}_y\) is a maximal element of \(\mathbf{N}\) . Show that \(xx, x(x(xx))\) are primitive whereas \(x, (xx)(xx)\) are not.
\item
  Let \(\mathbf{P}\) be the set of primitive elements of \(\mathbf{M}_x\) . Show that, if \(y, z \in \mathbf{P}, y \neq z\) , then \(\mathbf{M}_y \cap \mathbf{M}_z = \varnothing\) (use Exercise 6).
\item
  Let \(\operatorname{Mo}(\mathbf{P})\) be the free monoid constructed on \(\mathbf{P}\) . If \(\mathbf{M}_x\) is given the monoid structure defined in (b), the injection \(\mathbf{P} \to \mathbf{M}_x\) can be extended to a homomorphism \(p: \operatorname{Mo}(\mathbf{P}) \to \mathbf{M}_x\) . Show that \(p\) is an isomorphism.
\end{enumerate}

(If \(z \in \mathbf{M}_x\) , show by induction on \(l(z)\) that \(z \in \operatorname{Im}(p)\) . On the other hand, if \(y_1, \ldots, y_n, z_1, \ldots, z_m \in \mathbb{P}\) are such that \(p(y_1 \ldots y_n) = p(z_1 \ldots z_m)\) , then \(\mathbf{M}_{y_1} \cap \mathbf{M}_{z_1} \neq \varnothing\) , whence \(y_1 = z_1\) and \(f_{y_1}(p(y_2 \ldots y_n)) = f_{y_1}(p(z_2 \ldots z_m))\) ; whence \(p(y_2 \ldots y_n) = p(z_2 \ldots z_m)\) and it follows that \(y_1 \ldots y_n = z_1 \ldots z_m\) arguing by induction on \(\sup(n, m)\) .)

\begin{enumerate}
\def\labelenumi{\arabic{enumi}.}
\setcounter{enumi}{7}
\tightlist
\item
  Let X be a set; for every integer \(q \geqslant 0\) let \(\mathrm{Mo}(\mathbf{X})_{q}\) be the set of elements of \(\mathrm{Mo}(\mathbf{X})\) of length q and let \(\mathrm{Mo}^{(q)}(\mathbf{X})\) be the union of \(\mathrm{Mo}(\mathbf{X})_{nq}, n \in \mathbb{N}\) . The injection \(\mathrm{Mo}(\mathbf{X})_{q} \to \mathrm{Mo}^{(q)}(\mathbf{X})\) extends to a homomorphism
\end{enumerate}

\[
\operatorname{Mo} (\operatorname{Mo} (\mathrm{X}) _ {q}) \rightarrow \operatorname{Mo} ^ {(q)} (\mathrm{X});
\]

show that this homomorphism is bijective if \(q \geqslant 1\) .

\begin{enumerate}
\def\labelenumi{\arabic{enumi}.}
\setcounter{enumi}{8}
\item
  Let \(x, y \in \mathbf{X}\) with \(x \neq y\) . Show that the submonoid of \(\operatorname{Mo}(\mathbf{X})\) generated by \(\{x, xy, yx\}\) is not isomorphic to a free monoid.
\item
  Show that the group defined by the presentation
\end{enumerate}

\[
\langle x, y; x y ^ {2} = y ^ {3} x, y x ^ {2} = x ^ {3} y \rangle
\]

reduces to the identity element.

(The first relation implies \(x^{2}y^{8}x^{-2} = y^{18}\) and \(x^{3}y^{8}x^{-3} = y^{27}\) ; use the second relation to deduce that \(y^{18} = y^{27}\) , whence \(y^{9} = e\) ; the fact that \(y^{2}\) is conjugate to \(y^{3}\) then implies \(y = e\) , whence \(x = e\) .)

\begin{enumerate}
\def\labelenumi{\arabic{enumi}.}
\setcounter{enumi}{10}
\tightlist
\item
  The group defined by the presentation \(\langle x,y;x^{2},y^{3},(xy)^{2}\rangle\) is isomorphic to \(S_{3}\) .
\end{enumerate}

¶ 12. The group defined by the presentation \(\langle x, y; x^3, y^2, (xy)^3 \rangle\) is isomorphic to \(\mathfrak{A}_4\) . (Show that the conjugates of \(y\) form, with the identity element, a normal subgroup of order \(\leqslant 4\) and index \(\leqslant 3\) .)

\begin{enumerate}
\def\labelenumi{\arabic{enumi}.}
\setcounter{enumi}{12}
\tightlist
\item
  Let \(G\) be a group and \(X\) a subset of \(G\) disjoint from \(X^{-1}\) . Let \(S = X \cup X^{-1}\) . Show the equivalence of the two following properties:
\end{enumerate}

\begin{enumerate}
\def\labelenumi{(\roman{enumi})}
\item
  X is a free family in G.
\item
  No product \(s_1 \ldots s_n\) , with \(n \geqslant 1\) , \(s_i \in S\) and \(s_i s_{i+1} \neq e\) for \(1 \leqslant i \leqslant n-1\) , is equal to \(e\) .
\end{enumerate}

\begin{enumerate}
\def\labelenumi{\arabic{enumi}.}
\setcounter{enumi}{13}
\tightlist
\item
  A group G is called free if it possesses a basic family and hence is isomorphic to a free group F(X) constructed on a set X.
\end{enumerate}

\begin{enumerate}
\def\labelenumi{(\alph{enumi})}
\tightlist
\item
  Show that, if \((t_i)_{i\in I}\) and \((t_j')_{j\in J}\) are two such families, then \(\operatorname{Card}(I) = \operatorname{Card}(J)\) . (When \(\operatorname{Card}(I)\) is infinite, show that \(\operatorname{Card}(J)\) is also and that both are equal to \(\operatorname{Card}(G)\) . When \(\operatorname{Card}(I)\) is an integer \(d\) , show that the number of subgroups of index 2 in \(G\) is \(2^d - 1\) .)
\end{enumerate}

The cardinal of a basic family of G is called the rank of G.

\begin{enumerate}
\def\labelenumi{(\alph{enumi})}
\setcounter{enumi}{1}
\item
  A free group is of rank 0 (resp. 1) if and only if it reduces to \(\{e\}\) (resp. is isomorphic to \(\mathbf{Z}\) ).
\item
  Let G be a free group of finite rank d and \((x_{1},\ldots,x_{d})\) a generating family of d elements of G. Show that this family is basic. (Use Exercise 34 and § 5, Exercise 5.)
\item
  Show that a free group of rank \(\geqslant 2\) contains a free subgroup of given rank \(d\) for every cardinal \(d \leqslant \aleph_0\) . (Use Exercise 22.)
\end{enumerate}

¶ 15. Let G be a group with presentation (t, r), where \(\mathbf{t} = (t_{i})_{i \in I}\) and \(\mathbf{r} = (r_{j})_{j \in J}\) ; let \(\mathrm{F}((\mathrm{T}_{i})_{i \in I})\) denote the free group \(\mathrm{F}(I)\) and \(\mathbf{T} = (\mathrm{T}_{i})\) . On the other hand, let \(x = (x_{\alpha})_{\alpha \in A}\) be a generating family of G; write

\[
\mathrm{F} (\mathbf {A}) = \mathrm{F} ((\mathbf {X} _ {\alpha}) _ {\alpha \in A}),
\]

\(\mathbf{X} = (\mathbf{X}_{\alpha})\) and \(\mathbf{R}_{\mathbf{x}}\) the set of relators of \(\mathbf{x}\) ; it is a normal subgroup of \(\mathrm{F(A)}\) . For all \(i \in \mathbf{I}\) , let \(\phi_i(\mathbf{X})\) be an element of \(\mathrm{F(A)}\) such that \(t_i = \phi_i(\mathbf{x})\) in \(\mathbf{G}\) ; for all \(\alpha \in \mathbf{A}\) , let \(\psi_{\alpha}(\mathbf{T})\) be an element of \(\mathrm{F(I)}\) such that \(x_{\alpha} = \psi_{\alpha}(\mathbf{t})\) in \(\mathbf{G}\) . Show that \(\mathbf{R}_{\mathbf{x}}\) is the normal subgroup of \(\mathrm{F(A)}\) generated by the elements \(\mathbf{X}_{\alpha}^{-1} \psi_{\alpha}(\phi_i(\mathbf{X})_{i \in I}), \alpha \in \mathbf{A}\) and \(r_j(\phi_i(\mathbf{X})_{i \in I}) j \in J\) . (If \(\mathbf{R}_{\mathbf{x}}'\) denotes the normal subgroup of \(\mathrm{F(A)}\) generated by the elements in question, then \(\mathbf{R}_{\mathbf{x}}' \subset \mathbf{R}_{\mathbf{x}}\) ; on the other hand, show that the homomorphism \(\mathrm{F(I)} \to \mathrm{F(A)}\) defined by the \(\phi_i\) gives, when passing to the quotient, a homomorphism \(\mathrm{G} \to \mathrm{F(A)} / \mathrm{R}_{\mathbf{x}}'\) the inverse of the canonical homomorphism \(\mathrm{F(A)} / \mathrm{R}_{\mathbf{x}}' \to \mathrm{F(A)} / \mathrm{R}_{\mathbf{x}} = \mathrm{G}\) .)

\begin{enumerate}
\def\labelenumi{\arabic{enumi}.}
\setcounter{enumi}{15}
\tightlist
\item
  A group is called finitely generated (resp. finitely presented) if it admits a presentation \(\langle (t_i)_{i\in I}, (r_j)_{j\in J}\rangle\) where \(I\) is a finite set (resp. where \(I\) and \(J\) are finite sets).
\end{enumerate}

\begin{enumerate}
\def\labelenumi{(\alph{enumi})}
\item
  Let G be a finitely presented group and \(\mathbf{x} = (x_{\alpha})_{\alpha \in \mathbf{A}}\) a finite generating family of G. Show that there exists a finite family \(\mathbf{s} = (s_{k})_{k \in \mathbf{K}}\) of relators of the family x such that \((\mathbf{x}, \mathbf{s})\) is a presentation of G. (Use the preceding exercise.)
\item
  Show that a finite group is finitely presented.
\item
  Give an example of a finitely presented group with a subgroup which is not finitely generated.
\item
  If \(\mathbf{G}\) is a finitely presented group and \(\mathbf{H}\) a normal subgroup of \(\mathbf{G}\) , show the equivalence of the following properties:
\item
  G/H is finitely presented.
\end{enumerate}

\begin{enumerate}
\def\labelenumi{(\roman{enumi})}
\setcounter{enumi}{1}
\item
  There exists a finite subset X of G such that H is the normal subgroup of G generated by X.
\item
  If \((\mathbf{H}_i)_{i\in \mathbf{I}}\) is a right directed family of normal subgroup of \(\mathbf{G}\) whose union is equal to \(\mathbf{H}\) , there exists \(i\in \mathbf{I}\) such that \(\mathbf{H}_i = \mathbf{H}\) .
\end{enumerate}

(Use (a) to prove that (i) implies (ii).)

\begin{enumerate}
\def\labelenumi{(\alph{enumi})}
\setcounter{enumi}{4}
\item
  If G is a finitely presented group, show that \(\mathbf{G}/\mathbf{C}^{r}(\mathbf{G})\) is finitely presented for all \(r \geqslant 1\) . (Use §6, Exercise 15.)
\item
  Let \((\mathbf{G}_{\alpha})_{\alpha \in \mathbf{A}}\) be a family of groups and let \(\mathbf{G} = \prod_{\alpha \in \mathbf{A}} \mathbf{G}_{\alpha}\) be the product of the \(\mathbf{G}_{\alpha}\) . Show that \(\mathbf{G}\) is finitely generated (resp. finitely presented) if and only if each of the \(\mathbf{G}_{\alpha}\) is finitely generated (resp. finitely presented) and \(\mathbf{G}_{\alpha} = \{e\}\) for all but a finite number of \(\alpha\) .
\end{enumerate}

¶ 17. (a) Show that every subgroup of a finitely generated nilpotent group is finitely generated. (In the commutative case, argue by induction on the minimum number of generators of the group; then proceed by induction on the nilpotency class of the group.)

Give an example of a finitely generated solvable group containing a subgroup which is not finitely generated.

\begin{enumerate}
\def\labelenumi{(\alph{enumi})}
\setcounter{enumi}{1}
\tightlist
\item
  Show that every finitely generated nilpotent group is finitely presented. (Write the group in the form \(\mathrm{F}(\mathbf{X}) / \mathbb{R}\) , with \(\mathbf{X}\) finite, and choose \(r\) such that \(\mathbf{R} \supset \mathbf{C}^r (\mathrm{F}(\mathbf{X}))\) ; use (a) to show that \(\mathbb{R} / \mathbb{C}^r (\mathrm{F}(\mathbf{X}))\) is finitely generated; observe that \(\mathrm{F}(\mathbf{X}) / \mathbb{C}^r (\mathrm{F}(\mathbf{X}))\) is finitely presented, cf.~Exercise 16.)
\end{enumerate}

\begin{enumerate}
\def\labelenumi{\arabic{enumi}.}
\setcounter{enumi}{17}
\tightlist
\item
  Let S be a symmetric subset of a group G. Two elements \(g_{1}, g_{2}\) of G are called S-neighbours if \(g_{2}^{-1}g_{1}\) belongs to S. A sequence \((g_{1}, \ldots, g_{n})\) of elements of G is called an S-chain if \(g_{i}\) and \(g_{i+1}\) are S-neighbours for \(1 \leqslant i \leqslant n - 1\) ; the elements \(g_{1}\) and \(g_{n}\) are called respectively the beginning and the end of the chain.
\end{enumerate}

\begin{enumerate}
\def\labelenumi{(\alph{enumi})}
\item
  Let Y be a subset of G and \(R_{Y}\{a, b\}\) the relation ``there exists an S-chain in Y beginning at a and ending at \(b''\) . Show that \(R_{Y}\) is an equivalence relation on Y (cf.~Set Theory, II, §6, Exercise 10); the equivalence classes under this relation are called the connected S-components of Y. Y is called S-connected if it has at most one connected component.
\item
  Show that \(\mathbf{G}\) is S-connected if and only if \(\mathbf{S}\) generates \(\mathbf{G}\) .
\item
  Suppose that no element \(s \in S\) satisfies the relation \(s^2 = e\) ; choose a subset \(X\) of \(S\) such that \(S\) is the disjoint union of \(X\) and \(X^{-1}\) . Show the equivalence of the following conditions:
\item
  The family \(\mathbf{X}\) is free.
\end{enumerate}

\begin{enumerate}
\def\labelenumi{(\roman{enumi})}
\setcounter{enumi}{1}
\tightlist
\item
  There exists no S-chain \((g_{1},\ldots ,g_{n})\) in G consisting of distinct elements \(n\geqslant 3\) in number such that \(g_{1}\) and \(g_{n}\) are S-neighbours.
\end{enumerate}

(Use Exercise 13.)

¶ 19. Let X be a set, G = F(X) the free group constructed on X and S = X ∪ X \(^{-1}\) .

\begin{enumerate}
\def\labelenumi{(\alph{enumi})}
\item
  If \(g \in G\) , every sequence \((s_1, \ldots, s_n)\) of elements of \(S\) such that \(g = s_1 \ldots s_n\) and \(s_i s_{i+1} \neq e\) for \(1 \leqslant i < n\) is called a reduced decomposition of \(g\) . Show that every element of \(G\) admits one and only one reduced decomposition; the integer \(n\) is the length of \(g\) (cf.~Exercise 26).
\item
  Let Y be a subset of G containing e. Show the equivalence of the following conditions:
\end{enumerate}

\((b_{1})\) Y is S-connected (cf.~Exercise 18).

\((b_{2})\) For all \(g \in \mathbf{Y}\) , if \((s_1, \ldots, s_n)\) is the reduced decomposition of \(g\) , then \(s_1 \ldots s_i \in \mathbf{Y}\) for \(1 \leqslant i \leqslant n\) .

\begin{enumerate}
\def\labelenumi{(\alph{enumi})}
\setcounter{enumi}{2}
\tightlist
\item
  Let \(Y_{1}\) and \(Y_{2}\) be two non-empty S-connected subsets of G such that \(Y_{1} \cap Y_{2} = \varnothing\) . Show that there exists at most one ordered pair
\end{enumerate}

\[
\left(y _ {1}, y _ {2}\right) \in \mathbf {Y} _ {1} \times \mathbf {Y} _ {2}
\]

such that \(y_{1}\) and \(y_{2}\) are S-neighbours. There exists one if and only if \(\mathbf{Y}_1 \cup \mathbf{Y}_2\) is S-connected; \(\mathbf{Y}_1\) and \(\mathbf{Y}_2\) are then said to be neighbours.

\begin{enumerate}
\def\labelenumi{(\alph{enumi})}
\setcounter{enumi}{3}
\tightlist
\item
  Let \(Y_{1}, \ldots, Y_{n}\) be a sequence of disjoint non-empty S-connected subsets of G. Suppose that \(Y_{i}\) and \(Y_{i+1}\) are neighbours for \(1 \leqslant i < n\) . Show that, if \(n \geqslant 3\) , \(Y_{1}\) and \(Y_{n}\) are not neighbours.
\end{enumerate}

¶ 20. We preserve the notation of the preceding exercise. Let H be a subgroup of G.

\begin{enumerate}
\def\labelenumi{(\alph{enumi})}
\item
  Let \(R_{H}\) be the set of S-connected subsets T of G such that \(hT \cap T = \varnothing\) if \(h \in H, h \neq e\) . Show that \(R_{H}\) is inductive for the relation of inclusion.
\item
  Let Y be a maximal element of \(R_{H}\) . Show that \(G = \bigcup_{h \in H} hY\) , in other words Y is a system of representatives of the right cosets modulo H. (If \(G' = \bigcup hY\) is distinct from G, show that there exist \(x \in G'\) and \(y \in Y\) which are S-neighbours and deduce that \(Y \cup \{x\}\) belongs to \(R_{H}\) .)
\item
  Let \(S_{H}\) be the set of elements \(h \in H\) such that Y and hY are neighbours. It is a symmetric subset of H. Show that an element \(h \in H\) belongs to \(S_{H}\) if and only if \(h \neq e\) and there exist \(y, y' \in Y, s \in S\) with \(ys = hy'\) .
\end{enumerate}

Show that \(S_{H}\) generates H. (If \(H'\) is the subgroup of H generated by \(S_{H}\) .

prove that \(\bigcup_{h\in H'}hY\) is a connected S-component of G.) Obtain this result by means of Exercise 14 of §4.

\begin{enumerate}
\def\labelenumi{(\alph{enumi})}
\setcounter{enumi}{3}
\tightlist
\item
  Let \(X_{H}\) be a subset of \(S_{H}\) such that \(S_{H}\) is the disjoint union of \(X_{H}\) and \(X_{H}^{-1}\) (show that such a subset exists).
\end{enumerate}

Show that \(X_{H}\) is a basic family for H and in particular that H is a free group (``Nielsen-Schreier Theorem''). (Apply the criterion of Exercise 18(c) to H together with \(S_{H}\) , noting that two elements h, \(h'\) of H are \(S_{H}\) -neighbours if and only if hY and \(h'Y\) are neighbouring subsets of G. Use Exercise 19(d) to prove that \(S_{H}\) satisfies condition (ii) in Exercise 18(c).)

\begin{enumerate}
\def\labelenumi{(\alph{enumi})}
\setcounter{enumi}{4}
\tightlist
\item
  Let \(d = (\mathbf{G}:\mathbf{H}) = \operatorname{Card}(\mathbf{Y})\) . Suppose \(d\) is finite. Show that:
\end{enumerate}

\[
\begin{array}{l l} \operatorname{Card} (\mathbf {X} _ {\mathrm{H}}) = \operatorname{Card} (\mathbf {X}) = \operatorname{Card} (\mathbf {G}) & \text { if } \operatorname{Card} (\mathbf {X}) \text { is   infinite } \\ \operatorname{Card} (\mathbf {X} _ {\mathrm{H}}) = 1 + d (\operatorname{Card} (\mathbf {X}) - 1) & \text { if } \operatorname{Card} (\mathbf {X}) \text { is   finite }. \end{array}
\]

(Let T be a non-empty S-connected finite subset of G and T' (resp. T'') the set of ordered pairs of S-neighbouring elements whose first element (resp. both of whose elements) belongs (resp. belong) to T. Card(T') = 2Card(X)Card(T) and it should be shown by induction on Card(T) that

\[
\operatorname{Card} \left(\mathrm{T} ^ {\prime \prime}\right) = 2 \operatorname{Card} (\mathrm{T}) - 2.
\]

Apply the result to T = Y, noting that \(S_{H}\) is equipotent to \(T' - T''\) .

\begin{enumerate}
\def\labelenumi{\arabic{enumi}.}
\setcounter{enumi}{20}
\item
  Let H be a subgroup of a group G of finite index. Show that H is finitely presented if and only if G is. (It may be assumed that G is finitely generated, cf.~§4, Exercise 14. Write G in the form G = F/R with F free and finitely generated; the inverse image F' of H in F is free and finitely generated, cf.~Exercise 20, and H = F'/R. If R is generated (as a normal subgroup of F) by \((r_{j})_{j\in J}\) , show that it is generated (as a normal subgroup of F') by the \(yr_{j}y^{-1}\) , \(y\in Y\) , \(j\in J\) , where Y denotes a representative system in F of the right cosets modulo F'.)
\item
  \begin{enumerate}
  \def\labelenumii{(\alph{enumii})}
  \tightlist
  \item
    Let \((y_{n})_{n\in \mathbf{Z}}\) be a basic family of a group F. Let \(\phi\) be the automorphism of F which maps \(y_{n}\) to \(y_{n + 1}\) ; this automorphism defines an operation \(\tau\) of \(\mathbf{Z}\) on F; let \(\mathbf{E} = \mathbf{F}\times_{\tau}\mathbf{Z}\) be the corresponding semi-direct product. Show that the elements \(x = (e,1)\) and \(y = (y_0,0)\) form a basic family of E. (Let \(\mathrm{F}_{x,y}\) be the free group constructed on \(\{x,y\}\) and let \(f\) be the canonical homomorphism of \(\mathrm{F}_{x,y}\) into E. Show that there exists a homomorphism \(g\colon \mathbf{E}\to \mathbf{F}_{x,y}\) such that \(g((0,1)) = x\) and \(g((y_n,0)) = x^n yx^{-n}\) and \(f\) and \(g\) are inverses of one another.)
  \end{enumerate}
\end{enumerate}

\begin{enumerate}
\def\labelenumi{(\alph{enumi})}
\setcounter{enumi}{1}
\item
  Deduce that the normal subgroup of \(F_{x,y}\) generated by y is a free group with basic family \((x^{n}yx^{-n})_{n\in\mathbf{Z}}\) . Obtain this result by applying Exercise 20 to the representative system consisting of the \(x^{n}, n\in Z\) .
\item
  Extend the above to free groups of arbitrary rank.
\end{enumerate}

\begin{enumerate}
\def\labelenumi{\arabic{enumi}.}
\setcounter{enumi}{22}
\tightlist
\item
  Let \(\mathbf{F}\) be a free group with basic family \((x, y)\) , \(x \neq y\) . Show that the derived group of F has as basic family the family of commutators \((x^i, y^j)\) , \(i \in \mathbf{Z}\) , \(j \in \mathbf{Z}\) , \(i \neq 0\) , \(j \neq 0\) .
\end{enumerate}

(Use Exercise 20 or Exercise 32.)

\begin{enumerate}
\def\labelenumi{\arabic{enumi}.}
\setcounter{enumi}{23}
\tightlist
\item
  Let G be a group, \(S = G - \{e\}\) and \(F = F(S)\) . For all \(s \in S\) , let \(x_s\) denote the corresponding element of F. If \(s, t \in S\) , let \(r_{s,t}\) denote the element of F defined by:
\end{enumerate}

\[
r _ {s, t} = x _ {s} x _ {t} x _ {s t} ^ {- 1} \quad \text { if } s t \neq e \text { in } G
\]

\[
r _ {s, t} = x _ {s} x _ {t} \quad \text { if } s t = e \text { in } G.
\]

Let \(\phi\) be the homomorphism of F onto G which maps \(x_{s}\) to \(s\) and let R be the kernel of \(\phi\) .

Show that the family \((r_{s,t})\) for \(s, t \in \mathbf{S} \times \mathbf{S}\) , is a basic family of the group \(\mathbf{R}\) . (Apply Exercise 20 to the system of representatives of \(\mathbf{G}\) in \(\mathbf{F}\) consisting of \(e\) and the \(x_{s'}\) .)

\begin{enumerate}
\def\labelenumi{\arabic{enumi}.}
\setcounter{enumi}{24}
\tightlist
\item
  Let G be a group. Show the equivalence of the following properties:
\end{enumerate}

\begin{enumerate}
\def\labelenumi{(\alph{enumi})}
\item
  G is a free group.
\item
  For every group H and extension E of G by H there exists a section \(G \rightarrow E\) .
\end{enumerate}

(To prove that (b) implies (a), take E to be a free group and use Exercise 20.)

\begin{enumerate}
\def\labelenumi{\arabic{enumi}.}
\setcounter{enumi}{25}
\tightlist
\item
  Let \(\mathbf{X}\) be a set and \(g\) an element of the free group \(\mathrm{F}(\mathbf{X})\) . Let \(g = \prod_{\alpha=1}^{n} x_{\alpha}^{m(\alpha)}\) be the canonical decomposition of \(g\) as a product of powers of elements of \(\mathbf{X}\) (cf.~no. 5, Proposition 7). The integer \(l(g) = \sum_{\alpha} |m(\alpha)|\) is called the length of \(g\) .
\end{enumerate}

\begin{enumerate}
\def\labelenumi{(\alph{enumi})}
\item
  \(g\) is called cyclically reduced if either \(x_{1} \neq x_{n}\) or \(x_{1} = x_{n}\) and \(m(1)m(n) > 0\) . Show that every conjugacy class of \(F(X)\) contains one and only one cyclically reduced element; it is the element of minimum length of the class in question.
\item
  An element \(x \in \mathrm{F}(\mathbf{X})\) is called primitive if there does not exist an integer \(n \geqslant 2\) and an element \(g \in \mathrm{F}(\mathbf{X})\) such that \(x = g^n\) . Show that every element \(z \neq e\) of \(\mathrm{F}(\mathbf{X})\) can be written uniquely in the form \(x^n\) with \(n \geqslant 1\) and \(x\) primitive. (Reduce it to the case where \(z\) is cyclically reduced and observe that, if \(z = x^n\) , the element \(x\) is also cyclically reduced.)
\end{enumerate}

Show that the centralizer of z is the same as that of x; it is the cyclic subgroup generated by x.

\begin{enumerate}
\def\labelenumi{\arabic{enumi}.}
\setcounter{enumi}{26}
\tightlist
\item
  Let \(G = \times_{A} G_{i}\) be the sum of a family \((G_{i})_{i \in I}\) of groups amalgamated by a common subgroup A. For all \(i \in I\) , choose a subset \(P_{i}\) of \(G_{i}\) satisfying condition (A) of no. 3.
\end{enumerate}

\begin{enumerate}
\def\labelenumi{(\alph{enumi})}
\tightlist
\item
  Let \(x \in G\) and let \((a; i_1, \ldots, i_n; p_1, \ldots, p_n)\) be a reduced decomposition of \(x\) ; then
\end{enumerate}

\[
x = a \prod_ {\alpha = 1} ^ {n} p _ {\alpha}, \quad \text { with } \quad p _ {\alpha} \in \mathrm{P} _ {i _ {\alpha}} - \{e _ {i _ {\alpha}} \}, i _ {\alpha} \neq i _ {\alpha + 1}.
\]

Show that, if \(i_1 \neq i_n\) , the order of \(x\) is infinite. Show that, if \(i_1 = i_n\) , \(x\) is conjugate to an element with a decomposition of length \(\leqslant n - 1\) .

\begin{enumerate}
\def\labelenumi{(\alph{enumi})}
\setcounter{enumi}{1}
\tightlist
\item
  Deduce that every element of G of finite order is conjugate to an element of one of the \(G_{i}\) . In particular, if the \(G_{i}\) have no element of finite order except e, the same is true of G.
\end{enumerate}

\begin{enumerate}
\def\labelenumi{\arabic{enumi}.}
\setcounter{enumi}{27}
\tightlist
\item
  We preserve the notation of the preceding exercise.
\end{enumerate}

\begin{enumerate}
\def\labelenumi{(\alph{enumi})}
\item
  Let \(\mathbf{N}\) be a subgroup of \(\mathbf{A}\) . Suppose that, for all \(i \in \mathbf{I}\) , the group \(\mathbf{N}\) is normal in \(\mathbf{G}_i\) . Show that \(\mathbf{N}\) is normal in \(\mathbf{G} = \times_{\mathbf{A}} \mathbf{G}_i\) and that the canonical homomorphism of \(\times_{\mathbf{A/N}} \mathbf{G}_i / \mathbf{N}\) into \(\mathbf{G}/\mathbf{N}\) is an isomorphism.
\item
  For all \(i \in \mathbf{I}\) , let \(\mathbf{H}_i\) be a subgroup of \(\mathbf{G}_i\) and let \(\mathbf{B}\) be a subgroup of \(\mathbf{A}\) . Suppose that \(\mathbf{H}_i \cap \mathbf{A} = \mathbf{B}\) for all \(i\) . Show that the canonical homomorphism of \(\star_B \mathbf{H}_i\) into \(\mathbf{G}\) is injective; its image is the subgroup of \(\mathbf{G}\) generated by the \(\mathbf{H}_i\) .
\end{enumerate}

¶ 29. Let \(G = G_1 *_A G_2\) be the sum of two groups \(G_1\) and \(G_2\) amalgamated by a subgroup A.

\begin{enumerate}
\def\labelenumi{(\alph{enumi})}
\item
  Show that G is finitely generated if \(G_{1}\) and \(G_{2}\) are finitely generated.
\item
  Suppose \(\mathbf{G}_1\) and \(\mathbf{G}_2\) are finitely presented. Show the equivalence of the two following properties:
\end{enumerate}

\begin{enumerate}
\def\labelenumi{(\arabic{enumi})}
\item
  G is finitely presented.
\item
  A is finitely generated.
\end{enumerate}

(Show first that (1) implies (2). On the other hand, if A is not finitely generated, there exists a right directed family \((\mathbf{A}_{i})_{i\in I}\) of subgroups of A with \(\bigcup_{i\in I}\mathbf{A}_{i}=\mathbf{A}\) and \(A_{i}\neq A\) for all i. Let \(H_{i}\) (resp. H) be the kernel of the canonical homomorphism of \(G_{1}*G_{2}\) onto \(G_{1}*_{A_{i}}G_{2}\) (resp. onto G). Show that H is the union of the \(H_{i}\) and that \(H_{i}\neq H\) for all i; then apply Exercise 16(d).

\begin{enumerate}
\def\labelenumi{(\alph{enumi})}
\setcounter{enumi}{2}
\tightlist
\item
  Deduce an example of a finitely generated group which is not finitely presented. (Take \(G_{1}\) and \(G_{2}\) to be free groups of rank 2 and A a free group of finite rank, cf.~Exercise 22.)
\end{enumerate}

\begin{enumerate}
\def\labelenumi{\arabic{enumi}.}
\setcounter{enumi}{29}
\tightlist
\item
  Let A and B be two groups and G = A * B their free product.
\end{enumerate}

\begin{enumerate}
\def\labelenumi{(\alph{enumi})}
\tightlist
\item
  Let \(\mathbf{Z}\) be an element of \(\mathbf{A} \backslash \mathbf{G} / \mathbf{A}\) distinct from \(\mathbf{A}\) . Show that \(\mathbf{Z}\) contains one and only one element \(z\) of the form
\end{enumerate}

\[
b _ {1} a _ {1} b _ {2} a _ {2} \dots b _ {n - 1} a _ {n - 1} b _ {n},
\]

with \(n \geqslant 1\) , \(a_i \in \mathbf{A} - \{e\}\) , \(b_i \in \mathbf{B} - \{e\}\) . Show that every element of \(\mathbf{Z}\) can be written uniquely in the form \(aza'\) with \(a, a' \in \mathbf{A}\) .

\begin{enumerate}
\def\labelenumi{(\alph{enumi})}
\setcounter{enumi}{1}
\tightlist
\item
  Let \(x \in \mathbf{G} - \mathbf{A}\) . Let \(f_x\) be the homomorphism of \(\mathbf{A} * \mathbf{A}\) into \(\mathbf{G}\) whose restriction to the first (resp. second) factor of \(\mathbf{A} * \mathbf{A}\) is the identity (resp. the mapping \(a \mapsto xax^{-1}\) ). Show that \(f_x\) is injective. (Write \(x\) in the form \(aza'\) as above and reduce it thus to the case \(x = z\) ; use the uniqueness of reduced decompositions in \(\mathbf{G}\) .)
\end{enumerate}

In particular, \(\mathbf{A} \cap x\mathbf{A}x^{-1} = \{e\}\) and the normalizer of \(\mathbf{A}\) in \(\mathbf{G}\) is \(\mathbf{A}\) .

\begin{enumerate}
\def\labelenumi{(\alph{enumi})}
\setcounter{enumi}{2}
\tightlist
\item
  Let \(\mathbf{T}\) be an element of \(\mathbf{A} \backslash \mathbf{G} / \mathbf{B}\) . Show that \(\mathbf{T}\) contains one and only one element \(t\) of the form
\end{enumerate}

\[
b _ {1} a _ {1} b _ {2} a _ {2} \dots b _ {n - 1} a _ {n - 1},
\]

with \(n \geqslant 1\) , \(a_i \in \mathbf{A} - \{e\}\) , \(b_i \in \mathbf{B} - \{e\}\) . Show that every element of \(\mathbf{T}\) can be written uniquely in the form atb with \(a \in \mathbf{A}\) , \(b \in \mathbf{B}\) .

\begin{enumerate}
\def\labelenumi{(\alph{enumi})}
\setcounter{enumi}{3}
\tightlist
\item
  Let \(x \in G\) and let \(h_x\) be the unique endomorphism of \(G\) such that \(h_x(a) = a\) if \(a \in A\) and \(h_x(b) = x b x^{-1}\) if \(b \in B\) . Show that \(h_x\) is injective. (Same method as in (b), using the decomposition \(x = a t b\) of (c).) In particular, \(A \cap x B x^{-1} = \{e\}\) .
\end{enumerate}

Give an example where \(h_{x}\) is not surjective.

\begin{enumerate}
\def\labelenumi{\arabic{enumi}.}
\setcounter{enumi}{30}
\tightlist
\item
  Let G be the free product of a family \((G_{i})_{i\in I}\) of groups. Let S be the set of subgroups of G which are conjugate to one of the \(G_{i}\) and \(\Theta\) the union of the elements of S.
\end{enumerate}

\begin{enumerate}
\def\labelenumi{(\alph{enumi})}
\item
  Let \(\mathbf{H}, \mathbf{H}' \in \mathbf{S}\) , with \(\mathbf{H} \neq \mathbf{H}'\) , and let \(x \in \mathbf{H} - \{e\}\) , \(x' \in \mathbf{H}' - \{e\}\) . Show that \(xx' \in \mathbf{G} - \Theta\) . (Reduce it to the case where \(\mathbf{H}\) is one of the \(\mathbf{G}_i\) ; write \(\mathbf{H}'\) in the form \(y\mathrm{G}_j y^{-1}\) and proceed as in the preceding exercise.)
\item
  Let \(\mathbf{C}\) be a subgroup of \(\mathbf{G}\) contained in \(\Theta\) . Show that there exists \(\mathbf{H} \in \mathbf{S}\) such that \(\mathbf{C} \subset \mathbf{H}\) .
\item
  Let C be a subgroup of G all of whose elements are of finite order. Show that C is contained in \(\Theta\) (cf.~Exercise 27); deduce that C is contained in a conjugate of one of the \(G_{i}\) .
\end{enumerate}

¶ 32. Let A and B be two groups and G = A * B their free product. Let X be the set of commutators \((a, b)\) with \(a \in A - \{e\}\) and \(b \in B - \{e\}\) and \(S = X \cup X^{-1}\) .

\begin{enumerate}
\def\labelenumi{(\alph{enumi})}
\item
  Show that \(\mathbf{X} \cap \mathbf{X}^{-1} = \varnothing\) and that, if \(x\) is an element of \(\mathbf{X}\) , there exists a unique ordered pair \(a, b\) in \((\mathbf{A} - \{e\}) \times (\mathbf{B} - \{e\})\) such that \((a, b) = x\) .
\item
  Let \(n \geqslant 1\) and let \(s_1, \ldots, s_n\) be a sequence of elements of \(S\) such that \(s_i s_{i+1} \neq e\) for \(1 \leqslant i < n\) ; let \(a_n \in A - \{e\}, b_n \in B - \{e\}, \varepsilon(n) = \pm 1\) be such that \(s_n = (a_n, b_n)^{\varepsilon(n)}\) . Let \(g = s_1 \ldots s_n\) . Show, by induction on \(n\) , that the reduced decomposition of \(g\) is of length \(\geqslant n + 3\) and terminates either with \(a_nb_n\) if \(\varepsilon(n) = 1\) or with \(b_na_n\) if \(\varepsilon(n) = -1\) .
\item
  Let R be the kernel of the canonical projection \(A*B \rightarrow A \times B\) corresponding (no. 3) to the canonical injections \(A \rightarrow A \times B\) and \(B \rightarrow A \times B\) . Show that R is a free group with basic family X. (Use (b) to prove that X is free; if \(R_{X}\) is the subgroup of G generated by X, show that \(R_{X}\) is normal; deduce that it coincides with R.)
\end{enumerate}

\begin{enumerate}
\def\labelenumi{\arabic{enumi}.}
\setcounter{enumi}{32}
\tightlist
\item
  Let \(E = G *_{A} G'\) be the sum of two groups amalgamated by a common subgroup A. Let P (resp. \(P'\) ) be a subset of G (resp. \(G'\) ) satisfying condition (A) of no. 3.
\end{enumerate}

Let \(n\) be an integer \(\geqslant 1\) . Let \(L_n\) denote the set of \(x \in E\) whose reduced decomposition is of length \(\leqslant 2n - 1\) . Let \(\mathbf{M}_n\) denote the set of \(x \in \mathbf{E}\) of the form \(ap_1p_1'p_2p_2'\ldots p_np_n'\) , with \(a \in \mathbf{A}, p_i \in \mathbf{P} - \{e\}, p_i' \in \mathbf{P}' - \{e\}\) ; similarly, let \(\mathbf{M}_n'\) denote the set of elements of the form \(ap_1'p_1p_2'p_2\ldots p_np_n'\) , where \(a, p_i, p_i'\) satisfy the same conditions.

\begin{enumerate}
\def\labelenumi{(\alph{enumi})}
\item
  Show that \(L_{n}\) , \(M_{n}\) and \(M_{n}^{\prime}\) are disjoint and that their union is the set of elements of E whose reduced decomposition is of length \(\leqslant2n\) .
\item
  Construct a bijection \(\varepsilon\) of \(\mathbf{M}_n\) onto \(\mathbf{M}_n'\) such that \(\varepsilon(ax) = a\varepsilon(x)\) if \(a \in \mathbf{A}\) , \(x \in \mathbf{M}_n\) .
\item
  Let \(\mathbf{X}_n = \mathbf{L}_n \cup \mathbf{M}_n\) and \(\mathbf{X}_n' = \mathbf{L}_n \cup \mathbf{M}_n'\) ; show that \(\mathbf{G} \cdot \mathbf{X}_n = \mathbf{X}_n\) and \(\mathbf{G}' \cdot \mathbf{X}_n' = \mathbf{X}_n'\) .
\item
  \(\varepsilon\) extends to a bijection of \(\mathbf{X}_n\) onto \(\mathbf{X}_n'\) by setting \(\varepsilon(x) = x\) if \(x \in \mathbf{L}_n\) . Show that E can operate on \(\mathbf{X}_n\) so that \(g(x) = gx\) if \(g \in G\) , \(x \in \mathbf{X}_n\) and \(g'(x) = \varepsilon^{-1}(g' \varepsilon (x))\) if \(g' \in G'\) , \(x \in \mathbf{X}_n\) . Show that, if \(g \in \mathbf{X}_n\) , then \(g(e) = g\) .
\item
  Let \(\tau_{n}\colon\mathbf{E}\to\mathfrak{S}_{\mathbf{X}_{n}}\) be the homomorphism defined by the action of \(\mathbf{E}\) on \(\mathbf{X}_{n}\) described above and let \(\mathbf{R}_{n}\) be the kernel of \(\tau_{n}\) . Show that \(\mathbf{R}_{n}\cap\mathbf{X}_{n}=\{e\}\) ; deduce that the intersection of the \(\mathbf{R}_{n}\) reduces to \(e\) .
\item
  When G and G' are finite, show that E is residually finite (cf.~§ 5, Exercise 5). (Note that the R \(_{n}\) are then subgroups of finite index.)
\end{enumerate}

¶ 34. (a) Let \((\mathrm{H}_{i})\) be a right directed family of normal subgroups of a group G. Suppose that \(\bigcap_{i} \mathrm{H}_{i} = \{e\}\) . Show that, for every group \(\mathbf{G}'\) , the intersection of the kernels of the canonical homomorphisms \(\mathbf{G}' * \mathbf{G} \to \mathbf{G}' * (\mathbf{G}/\mathbf{H}_{i})\) is equal to \(\{e\}\) .

\begin{enumerate}
\def\labelenumi{(\alph{enumi})}
\setcounter{enumi}{1}
\item
  Let G be the free product of a family \((G_{\alpha})\) of groups. Show that, if all the \(G_{\alpha}\) are residually finite (cf.~§ 5, Exercise 5), so is G. (Reduce it first to the case where the family is finite, then to the case where it has two elements; by (a) it may be assumed that the \(G_{\alpha}\) are finite, then apply the preceding exercise.)
\item
  Deduce that every free group is residually finite.
\end{enumerate}

\begin{enumerate}
\def\labelenumi{\arabic{enumi}.}
\setcounter{enumi}{34}
\item
  *Let G (resp. G') be the additive group of fractions of the form \(a/b\) with \(a \in \mathbf{Z}\) , \(b \in \mathbf{Z}\) and \(b\) not divisible by 2 (resp. by 3). Let E = G * \(_{\mathbf{Z}}\) G' be the sum of G and G' amalgamated by their common subgroup Z. Show that G and G' are residually finite, but that E is not. (Prove that the only subgroup of E of finite index is E.)*
\item
  Let \((\mathbf{G}_1, \ldots, \mathbf{G}_n)\) and \((\mathbf{A}_1, \ldots, \mathbf{A}_{n-1})\) be two sequences of groups; for each \(i\) , let \(\phi_i: \mathbf{A}_i \to \mathbf{G}_i\) and \(\psi_i: \mathbf{A}_i \to \mathbf{G}_{i+1}\) be two injective homomorphisms; by these homomorphisms \(\mathbf{A}_i\) can be identified both with a subgroup of \(\mathbf{G}_i\) and with a subgroup of \(\mathbf{G}_{i+1}\) .
\end{enumerate}

\begin{enumerate}
\def\labelenumi{(\alph{enumi})}
\item
  Let \(H_{1}=G_{1}, H_{2}=H_{1} *_{A_{1}} G_{2}, \ldots, H_{n}=H_{n-1} *_{A_{n-1}} G_{n}\) . The group \(H_{n}\) is called the sum of the groups \((G_{i})\) amalgamated by the subgroups \((A_{i})\) ; it is denoted by \(G_{1} *_{A_{1}} G_{2} *_{A_{2}} \cdots *_{A_{n-1}} G_{n}\) . If T is an arbitrary group, define a bijection between the homomorphisms of \(H_{n}\) into T and the families \((t_{1},\ldots ,t_{n})\) of homomorphisms \(t_i\colon \mathbf{G}_i\to \mathbf{T}\) such that \(t_i\circ \phi_i = t_{i + 1}\circ \psi_i\) for \(1\leqslant i\leqslant n - 1\)
\item
  Each \(G_{i}\) is identified with the corresponding subgroup of \(H_{n}\) . Show that \(G_{i} \cap G_{i+1} = A_{i}\) and that \(G_{i} \cap G_{i+2} = A_{i} \cap A_{i+1}\) (the latter intersection being taken in \(G_{i+1}\) ).
\end{enumerate}

¶ 37. Let G be a group and S a finite symmetric subset of G generating G. Let \(\mathfrak{T}\) denote the set of finite subsets of G, ordered by inclusion; it is a directed set.

\begin{enumerate}
\def\labelenumi{(\alph{enumi})}
\item
  Let \(\mathbf{T} \in \mathfrak{T}\) and let \(\mathbf{B}_{\mathbf{T}}\) be the set of connected S-components of \(\mathbf{G} - \mathbf{T}\) (cf.~Exercise 18); let \(\mathbf{B}_{\mathbf{T}}^{\infty} = \mathbf{B}_{\mathbf{T}} - (\mathfrak{T} \cap \mathbf{B}_{\mathbf{T}})\) . Show that \(\mathbf{B}_{\mathbf{T}}\) is finite (note that, if \(\mathbf{T} \neq \varnothing\) and \(\mathbf{Z} \in \mathbf{B}_{\mathbf{T}}\) , there exists \(t \in \mathbf{T}\) and \(z \in \mathbf{Z}\) which are S-neighbours).
\item
  Let \(\mathbf{T}, \mathbf{T}' \in \mathfrak{T}\) with \(\mathbf{T} \subset \mathbf{T}'\) . If \(Z' \in \mathbf{B}_{\mathbf{T}'}\) , show that there exists a unique element \(Z \in \mathbf{B}_{\mathbf{T}}\) such that \(Z' \subset Z\) ; the mapping \(f_{\mathbf{T}\mathbf{T}'}: \mathbf{B}_{\mathbf{T}'} \to \mathbf{B}_{\mathbf{T}}\) such that \(f_{\mathbf{T}\mathbf{T}'}(\mathbf{Z}') = \mathbf{Z}\) maps \(\mathbf{B}_{\mathbf{T}}^{\infty}\) onto \(\mathbf{B}_{\mathbf{T}}^{\infty}\) . The inverse limit of the \(\mathbf{B}_{\mathbf{T}}\) , with \(\mathbf{T}\) running through \(\mathfrak{T}\) , is denoted by \(\mathbf{B}\) . Show that \(B = \varprojlim B_{\mathbf{T}}^{\infty}\) and that the image of \(B\) in \(B_{\mathbf{T}}\) is \(B_{\mathbf{T}}^{\infty}\) . The set \(B\) is called the set of ends of \(G\) ; it is non-empty if and only if \(G\) is infinite.
\item
  Let \(\mathbf{T}\) be a non-empty S-connected finite subset of \(\mathbf{G}\) and let \(\mathbf{Z} \in \mathbf{B}_{\mathbf{T}}^{\infty}\) . Show that there exist \(g \in \mathbf{G}\) such that \(g\mathbf{T} \cap \mathbf{T} = \varnothing\) and \(g\mathbf{T} \cap \mathbf{Z} \neq \varnothing\) and that then \(g\mathbf{T} \subset \mathbf{Z}\) . Show that there exists \(\mathbf{Z}_1 \in \mathbf{B}_{\mathbf{T}}\) such that \(g\mathbf{Z}_1\) contains \(\mathbf{T}\) and all the \(\mathbf{Z}' \in \mathbf{B}_{\mathbf{T}}\) such that \(\mathbf{Z}' \neq \mathbf{Z}\) (observe that \(\mathbf{T} \cup \bigcup_{\mathbf{Z}' \neq \mathbf{Z}} \mathbf{Z}'\) is S-connected and does not meet \(g\mathbf{T}\) ); deduce that, if \(\mathbf{Z}' \in \mathbf{B}_{\mathbf{T}}\) is different from \(\mathbf{Z}_1\) , then \(g\mathbf{Z}' \subset \mathbf{Z}\) . Show that, writing \(\mathbf{T}' = \mathbf{T} \cup g\mathbf{T}\) , the inverse image of \(\mathbf{Z}\) in \(\mathbf{B}_{\mathbf{T}}^{\infty}\) consists of at least \(n - 1\) elements, where \(n = \operatorname{Card}(\mathbf{B}_{\mathbf{T}}^{\infty})\) .
\item
  Let \(S'\) be a symmetric finite subset of \(G\) generating \(G\) and \(B'\) the set of ends of \(G\) relative to \(S'\) . Define a bijection of \(B\) onto \(B'\) (reduce it to the case where \(S \subset S'\) ). The cardinal of \(B\) is thus independent of the choice of \(S\) ; it is called the number of ends of \(G\) .
\item
  Show, using (c), that the number of ends of G is equal to 0, 1, 2 or \(2^{\aleph_0}\) .
\item
  Show that the number of ends of \(\mathbf{Z}\) is 2 and that that of \(\mathbf{Z}^n\) ( \(n \geqslant 2\) ) is 1.
\end{enumerate}

¶ 38. Let X be a finite set, F(X) the free group constructed on X and S = X ∪ X \(^{-1}\) .

\begin{enumerate}
\def\labelenumi{(\alph{enumi})}
\item
  Let \(\mathbf{S}_n\) be the set of products \(s_1 \ldots s_m\) , \(s_i \in \mathbf{S}\) , \(m \leqslant n\) . Show that every connected S-component (cf.~Exercise 18) of \(\mathrm{F}(\mathbf{X}) - \mathbf{S}_n\) contains one and only one element of the form \(s_1 \ldots s_{n+1}\) with \(s_i \in \mathbf{S}\) and \(s_i s_{i+1} \neq e\) for \(1 \leqslant i \leqslant n\) .
\item
  Deduce from (a) a bijection of the set B of ends of F(X) (cf.~Exercise 37) onto the set of infinite sequences \((s_{1},\ldots,s_{n},\ldots)\) of elements of S such that \(s_{i}s_{i+1}\neq e\) for all i. In particular, the number of ends of F(X) is equal to \(2^{\aleph_0}\) if \(\text{Card}(X)\geqslant2\) .
\end{enumerate}

\begin{enumerate}
\def\labelenumi{\arabic{enumi}.}
\setcounter{enumi}{38}
\tightlist
\item
  *Let G be a finitely generated group and F the set of symmetric finite subsets of G containing e, which generate G.
\end{enumerate}

\begin{enumerate}
\def\labelenumi{(\alph{enumi})}
\tightlist
\item
  If \(\mathbf{X} \in \mathfrak{F}\) and \(n \in \mathbf{N}\) , let \(d_n(\mathbf{X})\) denote the number of elements in \(\mathbf{G}\) of the form \(x_1 \ldots x_n\) with \(x_i \in \mathbf{X}\) . If \(\mathbf{X}, \mathbf{Y} \in \mathfrak{F}\) , show that there exists an integer \(a \geqslant 1\) such that
\end{enumerate}

\[
d _ {n} (\mathbf {X}) \leqslant d _ {a n} (\mathbf {Y}) \quad \text { for   all } n \in \mathbf {N}.
\]

Deduce that \(\limsup_{n\to \infty}\frac{\log(d_n(\mathbf{X}))}{\log(n)}\) does not depend on the choice of \(\mathbf{X}\) in \(\mathfrak{F}\) ; this limit (finite or infinite) is denoted by \(e(\mathbf{G})\) .

\begin{enumerate}
\def\labelenumi{(\alph{enumi})}
\setcounter{enumi}{1}
\item
  If \(\mathbf{X} \in \mathfrak{F}\) and \(\mathbf{G}\) is infinite, show that \(d_n(\mathbf{X}) < d_{n+1}(\mathbf{X})\) for all \(n\) . Deduce that \(e(\mathbf{G})\) is \(\geqslant 1\) .
\item
  Let H be a finitely generated group. Show that \(e(\mathbf{H}) \leqslant e(\mathbf{G})\) if H is isomorphic to a subgroup (resp. a quotient group) of G. If E is an extension of G by H, show that \(e(\mathbf{E}) \geqslant e(\mathbf{H}) + e(\mathbf{G})\) , with equality when \(E = G \times H\) .
\item
  Show that \(e(\mathbf{Z}^n) = n\) .
\item
  Let \(\mathbf{X} \in \mathfrak{F}\) . Consider the following property:
\item
  There exists a real number \(a > 1\) such that \(d_n(\mathbf{X}) \geqslant a^n\) for all sufficiently large \(n\) .
\end{enumerate}

Show that, if property (i) holds for one element of \(\mathfrak{F}\) , it holds for every element of \(\mathfrak{F}\) . \(\mathbf{G}\) is then said to be of exponential type; this implies \(e(\mathbf{G}) = +\infty\) .

\begin{enumerate}
\def\labelenumi{(\alph{enumi})}
\setcounter{enumi}{5}
\item
  Let \(G_{1}\) and \(G_{2}\) be two groups containing the same subgroup A and let \(H = G_{1} *_{A} G_{2}\) the corresponding amalgamated sum. Suppose \(G_{1}\) and \(G_{2}\) are finitely generated (hence also H is) and \(G_{i} \neq A\) for i = 1, 2. Show that H is of exponential type if at least one of the indices ( \(G_{1}:A\) ), ( \(G_{2}:A\) ) is \(\geqslant 3\) .
\item
  Show that a free group of rank \(\geqslant 2\) is of exponential type.*
\end{enumerate}

\begin{enumerate}
\def\labelenumi{\arabic{enumi}.}
\setcounter{enumi}{39}
\tightlist
\item
  *Let \(n\) be an integer \(\geqslant 2\) and \(T_n\) the group of upper triangular matrices of order \(n\) over \(Z\) with all the diagonal terms equal to 1. Let \(e(T_n)\) be the invariant of the group \(T_n\) defined in Exercise 39. Show that
\end{enumerate}

\[
e (\mathrm{T} _ {n}) \leqslant \sum_ {i = 1} ^ {n} i (n - i). *
\]

\BourbakiExercisesSection

\begin{enumerate}
\def\labelenumi{\arabic{enumi}.}
\item
  Determine all the ring structures on a set of \(n\) elements for \(2 \leqslant n \leqslant 7\) and also the ideals of these rings.
\item
  Let A be a ring.
\end{enumerate}

\begin{enumerate}
\def\labelenumi{(\alph{enumi})}
\item
  The mapping which associates with \(a \in A\) the left homothety \(x \mapsto ax\) is an isomorphism of A onto the ring of endomorphisms of the additive group of A which commute with right homotheties.
\item
  Let M be the set of sequences \((a, b, c, d)\) of four elements of A which may be represented in the form of a table or matrix \(\begin{pmatrix} a & b \\ c & d \end{pmatrix}\) . With every element \(\begin{pmatrix} a & b \\ c & d \end{pmatrix}\) of M we associate the mapping \((x, y) \mapsto (ax + by, cx + dy)\) of \(\mathbf{A} \times \mathbf{A}\)
\end{enumerate}

into itself. Show that a bijection is thus obtained of M onto the ring B of endomorphisms of the commutative group A × A which commute with the operations \((x, y) \mapsto (xr, yr)\) for all \(r \in A\) . Henceforth M is identified with the ring B. Calculate the product of two matrices.

\begin{enumerate}
\def\labelenumi{(\alph{enumi})}
\setcounter{enumi}{2}
\tightlist
\item
  Show that the subset \(\mathbf{T}\) of \(\mathbf{M}\) consisting of the matrices of the form \(\begin{pmatrix} a & b \\ 0 & c \end{pmatrix}\) is the subring of \(\mathbf{M}\) which leaves stable the subgroup \(\mathbf{A} \times 0\) of \(\mathbf{A} \times \mathbf{A}\) .
\end{enumerate}

The mappings \(\begin{pmatrix} a & b \\ 0 & c \end{pmatrix} \mapsto a\) and \(\begin{pmatrix} a & b \\ 0 & c \end{pmatrix} \mapsto c\) are homomorphisms of T into A; let \(a\) and \(b\) be their kernels. Show that \(a + b = T\) , that \(ab = \{0\}\) and that \(ba = b \cap a\) is \(\neq \{0\}\) if \(A \neq \{0\}\) (cf.~no. 9, Proposition 6).

\begin{enumerate}
\def\labelenumi{(\alph{enumi})}
\setcounter{enumi}{3}
\tightlist
\item
  If \(\mathbf{A} = \mathbf{Z} / 2\mathbf{Z}\) , \(\mathbf{T}\) is not commutative, but the multiplicative group of invertible elements of \(\mathbf{T}\) is commutative.
\end{enumerate}

\begin{enumerate}
\def\labelenumi{\arabic{enumi}.}
\setcounter{enumi}{2}
\item
  Let A be a ring and \(a \in A\) . If there exists one and only one \(a' \in A\) such that \(aa' = 1\) , a is invertible and \(a' = a^{-1}\) (show first that a is left cancellable, then consider the product \(aa'a\) ).
\item
  Let A be a pseudo-ring in which every additive subgroup of A is a left ideal of A.
\end{enumerate}

\begin{enumerate}
\def\labelenumi{(\alph{enumi})}
\item
  If \(\mathbf{A}\) is a ring, \(\mathbf{A}\) is canonically isomorphic to \(\mathbf{Z} / n\mathbf{Z}\) for some suitable integer \(n\) .
\item
  If A has no unit element and every non-zero element is cancellable, A is isomorphic to a pseudo-subring of Z.
\end{enumerate}

\begin{enumerate}
\def\labelenumi{\arabic{enumi}.}
\setcounter{enumi}{4}
\tightlist
\item
  Let \(\mathbf{M}\) be a monoid, \(\mathbf{Z}^{(\mathbf{M})}\) the free commutative group with basis \(\mathbf{M}\) and \(u\colon \mathbf{M}\to \mathbf{Z}^{(\mathbf{M})}\) the canonical mapping.
\end{enumerate}

\begin{enumerate}
\def\labelenumi{(\alph{enumi})}
\item
  Show that there exists a unique multiplication on \(\mathbf{Z}^{(\mathbf{M})}\) such that \(\mathbf{Z}^{(\mathbf{M})}\) is a ring and such that \(u\) is a monoid morphism of the monoid \(\mathbf{M}\) into the multiplicative monoid of the ring \(\mathbf{Z}^{(\mathbf{M})}\) .
\item
  Let A be a ring and \(v: \mathbf{M} \to \mathbf{A}\) a monoid morphism of the monoid M into the multiplicative monoid of A. There exists a unique ring morphism \(f: \mathbf{Z}^{(\mathbf{M})} \to \mathbf{A}\) such that
\end{enumerate}

\[
f \circ u = v.
\]

\begin{enumerate}
\def\labelenumi{\arabic{enumi}.}
\setcounter{enumi}{5}
\tightlist
\item
  Let A be a commutative ring and M a submonoid of the additive group of A. Let n be an integer \textgreater0 such that \(m^{n} = 0\) for \(m \in M\) .
\end{enumerate}

\begin{enumerate}
\def\labelenumi{(\alph{enumi})}
\item
  If \(n!\) is not a divisor of zero in \(\mathbf{A}\) , the product of a family of \(n\) elements of \(\mathbf{M}\) is zero.
\item
  Show that the conclusion in (a) is not necessarily true if \(n!\) is a divisor of zero in A.
\end{enumerate}

\begin{enumerate}
\def\labelenumi{\arabic{enumi}.}
\setcounter{enumi}{6}
\tightlist
\item
  For all \(n \in \mathbf{N}^*\) ,
\end{enumerate}

\[
n! = \sum_ {\nu = 0} ^ {n - 1} (- 1) ^ {\nu} \binom {n} {\nu} (n - \nu) ^ {n}
\]

(use no. 2, Proposition 2).

\begin{enumerate}
\def\labelenumi{\arabic{enumi}.}
\setcounter{enumi}{7}
\item
  In a ring, the right ideal generated by a left ideal is a two-sided ideal.
\item
  In a ring, the right annihilator of a right ideal is a two-sided ideal.
\item
  In a ring A, the two-sided ideal generated by the elements xy - yx, where x and y run through A, is the smallest of the two-sided ideals a such that A/a is commutative.
\end{enumerate}

¶ 11. In a ring A, a two-sided ideal a is said to be irreducible if there exists no ordered pair of two-sided ideals b, c distinct from a and such that a = b ∩ c.

\begin{enumerate}
\def\labelenumi{(\alph{enumi})}
\item
  Show that the intersection of all the irreducible two-sided ideals of A reduces to 0 (note that the set of two-sided ideals containing no element \(a \neq 0\) is inductive and apply Zorn's Lemma).
\item
  Deduce that every two-sided ideal \(\mathfrak{a}\) of \(\mathbf{A}\) is the intersection of all the irreducible ideals which contain it.
\end{enumerate}

\begin{enumerate}
\def\labelenumi{\arabic{enumi}.}
\setcounter{enumi}{11}
\tightlist
\item
  Let A be a ring and \((\mathfrak{a}_{i})_{i\in I}\) a finite family of left ideals of A such that the additive group of A is the direct sum of the additive groups \(a_{i}\) . If \(1=\sum_{i\in I}e_{i}(e_{i}\in\mathfrak{a}_{i})\) , then \(e_{i}^{2}=e_{i}, e_{i}e_{j}=0\) if \(i\neq j\) and \(a_{i}=Ae_{i}\) (write x=x.1 for all \(x\in A\) ). Conversely, if \((e_{i})_{i\in I}\) is a finite family of idempotents of A such that \(e_{i}e_{j}=0\) for \(i\neq j\) and \(1=\sum_{i}e_{i}\) , then A is the direct sum of the left ideals \(Ae_{i}\) . For the ideals \(Ae_{i}\) to be two-sided, it is necessary and sufficient that the \(e_{i}\) be in the centre of A.
\end{enumerate}

¶ 13. Let A be a ring and e an idempotent of A.

\begin{enumerate}
\def\labelenumi{(\alph{enumi})}
\item
  Show that the additive group of \(\mathbf{A}\) is the direct sum of the left ideal \(\mathfrak{a} = \mathbf{A}e\) and the left annihilator \(\mathfrak{b}\) of \(e\) (note that, for all \(x\in \mathbf{A}\) , \(x - xe\in \mathfrak{b}\) ).
\item
  Every right ideal \(\mathfrak{d}\) of \(\mathbf{A}\) is the direct sum of \(\mathfrak{d} \cap \mathfrak{a}\) and \(\mathfrak{d} \cap \mathfrak{b}\) .
\item
  If \(\mathbf{Ae} = e\mathbf{A}\) , \(\mathfrak{a}\) and \(\mathfrak{b}\) are two-sided ideals of \(\mathbf{A}\) and define a direct decomposition of \(\mathbf{A}\) .
\end{enumerate}

\begin{enumerate}
\def\labelenumi{\arabic{enumi}.}
\setcounter{enumi}{13}
\item
  Let A be a commutative ring in which there is only a finite number n of divisors of zero. Show that A has at most \((n+1)^{2}\) elements. (Let \(a_{1},\ldots,a_{n}\) be the divisors of zero. Show that the annihilator \(J_{1}\) of \(a_{1}\) has at most \(n+1\) elements and that the quotient ring A/ \(J_{1}\) has at most \(n+1\) elements.)
\item
  A ringoid is a set E with two laws of composition: (a) an associative multiplication \(xy\) ; (b) a law written additively, not everywhere defined (§ 1, no. 1) and satisfying the following conditions:
\end{enumerate}

\begin{enumerate}
\def\labelenumi{(\arabic{enumi})}
\item
  it is commutative (in other words, if \(x + y\) is defined, so is \(y + x\) and \(x + y = y + x\) ; \(x\) and \(y\) are then said to be addible);
\item
  if \(x\) and \(y\) are addible, for \(x + y\) and \(z\) to be addible it suffices that \(x\) and \(y\) on the one hand, and \(y\) and \(z\) on the other, be addible; then \(x\) and \(y + z\) are addible and \((x + y) + z = x + (y + z)\) ;
\item
  there exists an identity element 0;
\item
  if \(x\) and \(z\) on the one hand, \(y\) and \(z\) on the other, are addible and \(x + z = y + z\) , then \(x = y\) ;
\item
  if \(x, y\) are addible, so are \(xz\) and \(yz\) (resp. \(zx\) and \(zy\) ) and
\end{enumerate}

\[
(x + y) z = x z + y z
\]

(resp. \(z(x + y) = zx + zy)\) for all \(z \in \mathbf{E}\) .

Every ring is a ringoid.

\begin{enumerate}
\def\labelenumi{(\alph{enumi})}
\item
  Examine how the definitions and results of §8 and the above exercises extend to ringoids (a left ideal of a ringoid E is a subset \(\mathfrak{a}\) of E which is stable under addition and such that \(\mathbf{E}.\mathfrak{a}\subset \mathfrak{a}\) ).
\item
  Let G be a group with operators and f and g two endomorphisms of G. For the mapping \(x \mapsto f(x)g(x)\) to be an endomorphism of G, it is necessary and sufficient that every element of the subgroup \(f(\mathrm{G})\) be permutable with every element of \(g(\mathrm{G})\) ; if this endomorphism is then denoted by \(f + g\) and fg is the composite endomorphism \(x \mapsto f(g(x))\) , show that the set E of endomorphisms of G with these laws of composition is a ringoid with unit element; for E to be a ring, it is necessary and sufficient that G be commutative.
\end{enumerate}

For an element \(f \in E\) to be addible to all the elements of E, it is necessary and sufficient that \(f(G)\) be contained in the centre of G; the set N of these endomorphisms is a pseudo-ring called the kernel of the ringoid E.

An endomorphism \(f\) of \(G\) is called normal if it is permutable with all the inner automorphisms of \(G\) ; for every normal stable subgroup \(H\) of \(G, f(G)\) is then a normal stable subgroup of \(G\) . Show that the set \(D\) of normal endomorphisms of \(G\) forms a subringoid of \(E\) and that the kernel \(N\) is a two-sided ideal in \(D.†\)

\begin{enumerate}
\def\labelenumi{\arabic{enumi}.}
\setcounter{enumi}{15}
\tightlist
\item
  Give examples of ideals \(a, b, c\) in the ring \(\mathbf{Z}\) such that \(ab \neq a \cap b\) and \((a + b)(a + c) \neq a + bc\) .
\end{enumerate}

\BourbakiExercisesSection

\begin{enumerate}
\def\labelenumi{\arabic{enumi}.}
\item
  Which are the field structures amongst the ring structures determined in § 8, Exercise 1?
\item
  A finite ring in which every non-zero element is cancellable is a field (cf.~§ 2, Exercise 6).
\item
  Let G be a commutative group with operators which is simple (§ 4, no. 4, Definition 7). Show that the endomorphism ring A of G is a field.
\item
  Let K be a field. Show that there exists a smallest subfield F of K and that F is canonically isomorphic either to the field Q or to the field Z/pZ for some prime number p.~In the first case (resp. second case) K is said to be of characteristic 0 (resp. p).
\item
  Let K be a commutative field of characteristic \(\neq2\) ; let G be a subgroup of the additive group of K such that, if H denotes the set consisting of 0 and the inverses of the elements \(\neq0\) of G, H is also a subgroup of the additive group of K. Show that there exists an element \(a\in K\) and a subfield \(K'\) of K such that \(G=aK'\) (establish first that, if x and y are elements of G such that \(y\neq0\) , then \(x^{2}/y\in G\) ; deduce that, if x,y,z are elements of G such that \(z\neq0\) , then xy/z\ensuremath{\in}G).
\item
  Let K be a commutative field of characteristic \(\neq2\) ; let f be a mapping of K into K such that \(f(x+y)=f(x)+f(y)\) for all x and y and for all \(x\neq0\) \(f(x)f(1/x)=1\) . Show that f is an isomorphism of K onto a subfield of K (prove that \(f(x^{2})=(f(x))^{2}\) ).
\item
  Let A be a commutative ring.
\end{enumerate}

\begin{enumerate}
\def\labelenumi{(\alph{enumi})}
\item
  Every prime ideal is irreducible (§ 8, Exercise 11).
\item
  The set of prime ideals is inductive for the relations \(\subset\) and \(\supset\) .
\item
  Let \(\mathfrak{a}\) be an ideal of \(\mathbf{A}\) distinct from \(\mathbf{A}\) ; let \(\mathfrak{b}\) be the set of \(x \in \mathbf{A}\) such that there exists an integer \(n \geqslant 0\) with \(x^n \in \mathfrak{a}\) ( \(n\) depends on \(x\) ). Show that \(\mathfrak{b}\) is an ideal and is equal to the intersection of the prime ideals of \(\mathbf{A}\) containing \(\mathfrak{a}\) .
\end{enumerate}

¶ 8. A ring A is called a Boolean ring if each of its elements is idempotent (in other words, if \(x^2 = x\) for all \(x \in \mathbf{A}\) ).

\begin{enumerate}
\def\labelenumi{(\alph{enumi})}
\item
  In the set \(\mathfrak{P}(\mathbf{E})\) of subsets of a set \(\mathbf{E}\) , show that a Boolean ring structure is defined by setting \(\mathbf{AB} = \mathbf{A} \cap \mathbf{B}\) and \(\mathbf{A} + \mathbf{B} = (\mathbf{A} \cap \mathbb{C}\mathbf{B}) \cup (\mathbf{B} \cap \mathbb{C}\mathbf{A})\) . This ring is isomorphic to the ring \(\mathbf{K}^{\mathrm{E}}\) of mappings of \(\mathbf{E}\) into the ring \(\mathbf{K} = \mathbf{Z}/(2)\) of integers modulo 2 (consider for each subset \(\mathbf{X} \subset \mathbf{E}\) its ``characteristic function'' \(\phi_{\mathbf{X}}\) such that \(\phi_{\mathbf{X}}(x) = 1\) if \(x \in \mathbf{X}\) , \(\phi_{\mathbf{X}}(x) = 0\) if \(x \notin \mathbf{X}\) ).
\item
  Every Boolean ring \(\mathbf{A}\) is commutative and such that \(x + x = 0\) for \(x \in \mathbf{A}\) (write \(x + x\) as an idempotent and then \(x + y\) as an idempotent).
\item
  If a Boolean ring \(\mathbf{A}\) contains no divisor of 0, it is reduced to 0 or is isomorphic to \(\mathbf{Z} / (2)\) (if \(x\) and \(y\) are any two elements of \(\mathbf{A}\) , show that \(xy(x + y) = 0\) ). Deduce that in a Boolean ring every prime ideal is maximal.
\item
  In a Boolean ring A, every ideal \(a \neq A\) is the intersection of the prime ideals containing a (apply Exercise 7(c)). Deduce that every irreducible ideal is maximal (in other words, that the notions of irreducible ideal, prime ideal and maximal ideal coincide in a Boolean ring).
\item
  Show that every Boolean ring is isomorphic to a subring of a product ring \(\mathbf{K}^{\mathbb{E}}\) , where \(\mathbf{K} = \mathbf{Z}/(2)\) (use \((d)\) ).
\item
  If \(p_{i}\) ( \(1 \leqslant i \leqslant n\) ) are n distinct maximal ideals in a Boolean ring A and \(a = \bigcap_{1 \leqslant i \leqslant n} p_{i}\) , show that the quotient ring A/a is isomorphic to the product ring \(K^{n}\) . Deduce that every finite Boolean ring is of the form \(K^{n}\) .
\item
  In a Boolean ring A, the relation \(xy = x\) is an order relation; if it is denoted by \(x \leqslant y\) , show that with this relation A is a distributive lattice (Set Theory, III, § 1, Exercise 16), has a least element \(\alpha\) and that, for every ordered pair \((x, y)\) of elements such that \(x \leqslant y\) , there exists an element \(d(x, y)\) such that \(\inf(x, d(x, y)) = \alpha\) , \(\sup(x, d(x, y)) = y\) . Conversely, if an ordered set A has these three properties, show that the laws of composition \(xy = \inf(x, y)\) and \(x + y = d(\sup(x, y), \inf(x, y))\) define a Boolean ring structure on A.
\end{enumerate}

\begin{enumerate}
\def\labelenumi{\arabic{enumi}.}
\setcounter{enumi}{8}
\tightlist
\item
  Let \(\mathbf{A}\) be a ring such that \(x^3 = x\) for all \(x \in \mathbf{A}\) . It is proposed to prove that \(\mathbf{A}\) is commutative.
\end{enumerate}

\begin{enumerate}
\def\labelenumi{(\alph{enumi})}
\item
  Show that \(6\mathbf{A} = \{0\}\) and that \(2\mathbf{A}\) and \(3\mathbf{A}\) are two-sided ideals such that \(2\mathbf{A} + 3\mathbf{A} = \mathbf{A}\) and \(2\mathbf{A} \cap 3\mathbf{A} = \{0\}\) . Deduce that it can be assumed that either \(2\mathbf{A} = \{0\}\) or \(3\mathbf{A} = \{0\}\) .
\item
  If 2A = \{0\}, calculate \((1 + x)^{3}\) , deduce that \(x^{2} = x\) for all \(x \in A\) and conclude by means of Exercise 8.
\item
  If \(3\mathbf{A} = \{0\}\) , calculate \((x + y)^3\) and \((x - y)^3\) , show that \(x^2y + xyx + yx^2 = 0\) and left multiply by \(x\) ; deduce that \(xy - yx = 0\) .
\item
  Let A be a ring such that \(3A = \{0\}\) and \(x^{3} = x\) for all \(x \in A\) ; define a set I and an injective homomorphism of A into \((\mathbf{Z}/3\mathbf{Z})^{\mathrm{I}}\) (same method as in Exercise 8).
\end{enumerate}

\begin{enumerate}
\def\labelenumi{\arabic{enumi}.}
\setcounter{enumi}{9}
\tightlist
\item
  Let \(\mathbf{A}\) be a commutative ring and \(\mathfrak{U}\) the intersection of its maximal ideals.
\end{enumerate}

\begin{enumerate}
\def\labelenumi{(\alph{enumi})}
\item
  Show that the canonical homomorphism \(\mathbf{A}^* \to (\mathbf{A} / \mathfrak{U})^*\) is surjective and that its kernel is \(1 + \mathfrak{U}\) .
\item
  Suppose that the set M of maximal ideals of A is finite and that A* is finite. Deduce that A is finite (apply Proposition 8 of §8, no. 10 to A/U).
\item
  Deduce that the set of prime numbers is infinite (note that \(\mathbf{Z}^{*} = \{1, -1\}\) ).
\item
  Suppose that A* is finite and that A/m is finite for every maximal ideal m. Can it be deduced that A is finite, that A is countable?
\end{enumerate}

\begin{enumerate}
\def\labelenumi{\arabic{enumi}.}
\setcounter{enumi}{10}
\tightlist
\item
  Let P be the set of prime numbers and A the product ring of the fields \(\mathbf{Z} / p\mathbf{Z}\) , \(p \in \mathbb{P}\) . Let \(a\) be the subset of A consisting of the elements \((a_p)_{p \in \mathbb{P}}\) such that \(a_p \neq 0\) only for a finite number of indices \(p\) .
\end{enumerate}

\begin{enumerate}
\def\labelenumi{(\alph{enumi})}
\item
  a is an ideal of A.
\item
  Let \(\mathbf{B} = \mathbf{A} / a\) . For every integer \(n > 0\) and every \(b \neq 0\) in \(\mathbf{B}\) , there exists one and only one element \(b'\) of \(\mathbf{B}\) such that \(nb' = b\) (note that if \(p \in \mathbb{P}\) does not divide \(n\) , multiplication by \(n\) in \(\mathbf{Z} / p\mathbf{Z}\) is bijective). Deduce that \(\mathbf{B}\) contains a subfield isomorphic to \(\mathbf{Q}\) .
\end{enumerate}

\begin{enumerate}
\def\labelenumi{\arabic{enumi}.}
\setcounter{enumi}{11}
\item
  Let A be a commutative ring and \(a, b \in A\) . Show that the canonical image of ab in \(\mathrm{A}/(a - a^{2}b)\) is an idempotent. Give an example where this idempotent is distinct from 0 and 1.
\item
  Determine the endomorphisms of the ring Z and the ring \(Z \times Z\) . More generally, if I and J are finite sets, determine the homomorphisms of the ring \(Z^{1}\) into the ring \(Z^{J}\) .
\item
  Let A and B be two rings and f and g two homomorphisms of A into B. Is the set of \(x \in A\) such that \(f(x) = g(x)\) a subring of A, an ideal of A?
\end{enumerate}

¶ 15. Let A be a non-commutative ring with no divisors of 0. A is said to admit a field of left fractions if it is isomorphic to a subring B of a field K such that every element of K is of the form \(x^{-1}y\) , where \(x \in \mathbf{B}, y \in \mathbf{B}\) .

\begin{enumerate}
\def\labelenumi{(\alph{enumi})}
\tightlist
\item
  Let A' be the set of elements \(\neq0\) in A. For A to admit a field of left fractions, it is necessary that the following condition be fulfilled:
\end{enumerate}

\begin{enumerate}
\def\labelenumi{(\Alph{enumi})}
\setcounter{enumi}{6}
\tightlist
\item
  for all \(x \in \mathbf{A}\) , \(x' \in \mathbf{A}'\) , there exist \(u \in \mathbf{A}'\) and \(v \in \mathbf{A}\) such that \(ux = vx'\) .
\end{enumerate}

\begin{enumerate}
\def\labelenumi{(\alph{enumi})}
\setcounter{enumi}{1}
\tightlist
\item
  Suppose conversely that condition (G) is fulfilled. Show that in the set A × A' the relation R between \((x, x')\) and \((y, y')\) which states ``for every ordered pair \((u, v)\) of elements ≠0 such that \(ux' = vy'\) , \(ux = vy\) is an equivalence relation.
\end{enumerate}

Let \((x, x')\) , \((y, y')\) be two elements of \(A \times A'\) , \(\xi\) and \(\eta\) their respective classes (mod. R). For every ordered pair \((u, u') \in A \times A'\) such that \(u'x = uy'\) , show that the class (mod. R) of \((uy, u'y')\) depends only on the classes \(\xi\) and \(\eta\) ; if it is denoted by \(\xi\eta\) , a law of composition is defined on the set \(K = (A \times A')/R\) . If \(K'\) is the set of elements of K distinct from the class 0 of elements \((0, x')\) of \(A \times A'\) , \(K'\) , with the law induced by the above law, is a group.

For every element \(x \in \mathbf{A}\) , the elements \((x'x, x')\) , where \(x'\) runs through \(\mathbf{A}'\) , define an isomorphism of \(\mathbf{A}\) (with multiplication alone) onto a subring of \(\mathbf{K}\) . Identifying \(\mathbf{A}\) with its image under this isomorphism, the class (mod. R) of an ordered pair \((x, x') \in \mathbf{A} \times \mathbf{A}'\) is identified with the element \(x'^{-1}x\) .

This being so, if \(\xi = x'^{-1}x\) is an element of K and 1 denotes the unit element of \(\mathbf{K}'\) , then \(\xi + 1\) denotes the element \(x'^{-1}(x + x')\) , which does not depend on the representation of \(\xi\) in the form \(x'^{-1}x\) . We then write \(\xi + 0 = \xi\) and, for \(\eta \neq 0\) , \(\xi + \eta = \eta(\eta^{-1}\xi + 1)\) . Show that the addition and multiplication thus defined on K determine on that set a field structure which extends the ring structure on A; in other words, condition (G) is sufficient for A to admit a field of left fractions.

\begin{enumerate}
\def\labelenumi{\arabic{enumi}.}
\setcounter{enumi}{15}
\item
  Let A be a ring in which every non-zero element is cancellable. For A to admit a field of left fractions (Exercise 15), it is necessary and sufficient that in A the intersection of two left ideals distinct from 0 never reduce to 0.
\item
  Let \((\mathbf{K}_i)_{i\in I}\) be a family of fields. For every subset \(J\) of \(I\) , let \(e_J\) denote the element of the product \(\mathbf{K} = \prod_{i\in I}\mathbf{K}_i\) whose \(i\) -th component is equal to 0 if \(i\in J\) and to 1 if \(i\in I - J\) .
\end{enumerate}

\begin{enumerate}
\def\labelenumi{(\alph{enumi})}
\item
  Let \(x = (x_{i})\) be an element of \(\mathbf{K}\) and \(J(x)\) the set of elements \(i \in I\) such that \(x_{i} = 0\) . Show that there exists an invertible element \(u\) in \(\mathbf{K}\) such that \(x = ue_{J(x)} = e_{J(x)}u\) .
\item
  Let \(\mathfrak{a}\) be a left (resp. right) ideal of \(\mathbf{K}\) distinct from \(\mathbf{K}\) . Let \(\mathfrak{F}_{\mathfrak{a}}\) be the set of subsets \(\mathbf{J}\) of \(\mathbf{I}\) such that \(e_{J} \in \mathfrak{a}\) . Show that \(\mathfrak{F}_{\mathfrak{a}}\) has the following properties:
\end{enumerate}

\((b_{1})\) If \(\mathbf{J} \in \mathfrak{F}_{\mathfrak{a}}\) and \(\mathbf{J}' \supset \mathbf{J}\) , then \(\mathbf{J}' \in \mathfrak{F}_{\mathfrak{a}}\) .

\((b_{2})\) If \(\mathrm{J}_1, \mathrm{J}_2 \in \mathfrak{F}_{\mathfrak{a}}\) , then \(\mathrm{J}_1 \cap \mathrm{J}_2 \in \mathfrak{F}_{\mathfrak{a}}\) .

\[
\left(b _ {3}\right) \varnothing \notin \mathfrak {F} _ {\mathfrak {a}}.
\]

Show that \(x \in \mathfrak{a}\) is equivalent to \(J(x) \in \mathfrak{F}_{\mathfrak{a}}\) (use (a)). Deduce that \(\mathfrak{a}\) is a two-sided ideal.

\begin{enumerate}
\def\labelenumi{(\alph{enumi})}
\setcounter{enumi}{2}
\item
  A set of subsets of I is called a filter if it satisfies properties \((b_{1})\) , \((b_{2})\) , \((b_{3})\) above (*cf.~General Topology, I, §6*). Show that the mapping \(\mathfrak{a} \mapsto \mathfrak{F}_{\mathfrak{a}}\) is a strictly increasing bijection of the set of ideals of K distinct from K onto the set of filters of I. *Show that a is maximal if and only if \(\mathfrak{F}_{\mathfrak{a}}\) is an ultrafilter.*
\item
  *Let \(\mathfrak{F}\) be a non-trivial ultrafilter on the set \(P\) of prime numbers, let \(a\) be the corresponding ideal of the ring \(\mathbf{K} = \prod_{p \in P} \mathbf{Z}/p\mathbf{Z}\) and let \(k = K/a\) . Show that \(k\) is a field of characteristic \(0\)*
\end{enumerate}

¶ 18. (a) Let K be a field and a, b two non-permutable elements of K. Show that

\[
\begin{array}{l} (1) a = (b - (a - 1) ^ {- 1} b (a - 1)) (a ^ {- 1} b a - (a - 1) ^ {- 1} b (a - 1)) ^ {- 1} \\ (2) a = (1 - (a - 1) ^ {- 1} b ^ {- 1} (a - 1) b) (a ^ {- 1} b ^ {- 1} a b - (a - 1) ^ {- 1} b ^ {- 1} (a - 1) b) ^ {- 1}. \end{array}
\]

\begin{enumerate}
\def\labelenumi{(\alph{enumi})}
\setcounter{enumi}{1}
\item
  Let K be a non-commutative field and x an element not belonging to the centre of K. Show that K is generated by the set of conjugates \(axa^{-1}\) of x. (Let \(K_{1}\) be the subfield of K generated by the set of conjugates of x. Deduce from (a) that, if \(K_{1} \neq K\) and \(a \in K_{1}\) and \(b \notin K_{1}\) , then necessarily ab = ba. Deduce a contradiction by considering in \(K_{1}\) two non-permutable elements a, \(a'\) and an element \(b \notin K_{1}\) and noting that \(ba \notin K_{1}\) .)
\item
  Deduce from (b) that if Z is the centre of K and H is a subfield of K such that \(Z \subset H \subset K\) and \(aHa^{-1} = H\) for all \(a \neq 0\) in K, then H = Z or H = K (Cartan-Brauer-Hua Theorem).
\item
  Deduce from (a) that in a non-commutative field K, the set of commutators \(aba^{-1}b^{-1}\) of elements \(\neq0\) in K generates the field K.
\end{enumerate}

\begin{enumerate}
\def\labelenumi{\arabic{enumi}.}
\setcounter{enumi}{18}
\tightlist
\item
  Let A be a non-zero pseudo-ring such that, for all \(a \neq 0\) in A and all \(b \in \mathbf{A}\) , the equation \(ax + ya = b\) has solutions consisting of ordered pairs \((x, y)\) of elements of A; in other words, \(a\mathbf{A} + \mathbf{A}a = \mathbf{A}\) for \(a \neq 0\) in A.
\end{enumerate}

\begin{enumerate}
\def\labelenumi{(\alph{enumi})}
\tightlist
\item
  Show that A is the only non-zero two-sided ideal in A.
\end{enumerate}

\begin{enumerate}
\def\labelenumi{(\alph{enumi})}
\setcounter{enumi}{1}
\item
  Show that the relation \(a \neq 0\) implies \(a^2 \neq 0\) (note that if \(a^2 = 0\) , it would follow that \(aAa = \{0\}\) ).
\item
  Show that if \(ab = 0\) and \(a \neq 0\) (resp. \(b \neq 0\) ) then \(b = 0\) (resp. \(a = 0\) ). (Note that \((ba)^2 = 0\) and deduce that \(bAb = \{0\}\) .)
\end{enumerate}

\begin{enumerate}
\def\labelenumi{\arabic{enumi}.}
\setcounter{enumi}{19}
\item
  Let A be a non-zero pseudo-ring such that there exists \(a \in A\) for which, for all \(b \in A\) , one of the equations ax = b, xa = b admits a solution \(x \in A\) . Show that if A has no divisor of 0 other than 0, A admits a unit element.
\item
  Let A be a non-zero pseudo-ring in which, for all \(a \neq 0\) and all \(b \in A\) , one of the equations ax = b, xa = b admits a solution \(x \in A\) . Show that A is a field (use Exercises 19 and 20).
\end{enumerate}

\BourbakiExercisesSection

\begin{enumerate}
\def\labelenumi{\arabic{enumi}.}
\tightlist
\item
  Let X be the direct limit of a direct system \((\mathbf{X}_{i}, f_{ji})\) of sets relative to a right directed set I. Let \(F_{i}\) (resp. F) be the free group constructed on \(X_{i}\) (resp. on X); each \(f_{ji}\) extends to a homomorphism \(\phi_{ji}: F_{i} \to F_{j}\) .
\end{enumerate}

\begin{enumerate}
\def\labelenumi{(\alph{enumi})}
\item
  Show that \((\mathbf{F}_i,\phi_{ji})\) is a direct system of groups.
\item
  Let \(\varepsilon_{i}\) be the composite mapping \(\mathbf{X}_{i} \to \mathbf{F}_{i} \to \lim \mathbf{F}_{i}\) . Show that \(\varepsilon_{j} \circ \phi_{ji} = \varepsilon_{i}\) if \(j \geqslant i\) . Let \(\varepsilon\) be the mapping of \(\mathbf{X}\) into \(\varinjlim \mathbf{F}_{i}\) defined by the \(\varepsilon_{i}\) . Show that \(\varepsilon\) extends to an isomorphism of \(\mathbf{F}\) onto \(\varinjlim \mathbf{F}_{i}\) (so that \(\mathbf{F}(\varinjlim \mathbf{X}_{i})\) can be identified with \(\lim \mathbf{F}(\mathbf{X}_{i})\) ).
\item
  State and prove analogous results for free magmas, free monoids, free commutative groups, free commutative monoids.
\end{enumerate}

\begin{enumerate}
\def\labelenumi{\arabic{enumi}.}
\setcounter{enumi}{1}
\tightlist
\item
  Show that a direct limit of simple groups is either a simple group or a group with one element.
\end{enumerate}

\BourbakiStarSection{Historical Note}

(Numbers in brackets refer to the bibliography at the end of this Note.)

There are few notions in Mathematics more primitive than that of a law of composition: it seems inseparable from the first rudiments of calculations on natural numbers and measurable quantities. The most ancient documents that have come down to us on the Mathematics of the Egyptians and Babylonians show that they already had a complete system of rules for calculating with natural numbers \textgreater0, rational numbers \textgreater0, lengths and areas; although the preserved texts deal only with problems in which the given quantities have explicit numerical values,† they leave us in no doubt concerning the generality they attributed to the rules employed and show a quite remarkable technical ability in the manipulation of equations of the first and second degree ({[}1{]}, p.~179 et seq.). On the other hand there is not the slightest trace of a concern to justify the rules used nor even of precise definitions of the operations which occur: the latter as the former arise in a purely empirical manner.

Just such a concern, however, is already very clearly manifest in the works of the Greeks of the classical age; true, an axiomatic treatment is not found of the theory of natural numbers (such an axiomatization was only to appear at the end of the nineteenth century; see the Historical Note to Set Theory, IV); but in Euclid's Elements there are numerous passages giving formal proofs of rules of calculation which are just as intuitively ``obvious'' as those of calculation with integers (for example, the commutativity of the product of two rational numbers). The most remarkable proofs of this nature are those relating to the theory of magnitudes, the most original creation of Greek Mathematics (equivalent, as is known, to our theory of real numbers \textgreater0; see the Historical Note to General Topology, IV); here Euclid considers amongst other things the product of two ratios of magnitudes and shows that it is independent of the form of presentation of these ratios (the first example of a ``quotient'' of a law of composition by an equivalence relation in the sense of § 1, no. 6) and that they commute ({[}2{]}, Book V, Prop. 22--23).†

It must not however be concealed that this progress towards rigour is accompanied in Euclid by a stagnation and even in some ways by a retrogression as far as the technique of algebraic calculations is concerned. The overwhelming preponderance of Geometry (in the light of which the theory of magnitudes was obviously conceived) paralyses any independent development of algebraic notation: the elements entering into the calculations must always be ``represented'' geometrically; moreover the two laws of composition which occur are not defined on the same set (addition of ratios is not defined in a general way and the product of two lengths is not a length but an area); the result is a lack of flexibility making it almost impracticable to manipulate algebraic relations of degree greater than the second.

Only with the decline of classical Greek Mathematics does Diophantus return to the tradition of the ``logisticians'' or professional calculators, who had continued to apply such rules as they had inherited from the Egyptians and Babylonians: no longer encumbering himself with geometric representations of the ``numbers'' he considers, he is naturally led to the development of rules of abstract algebraic calculation; for example, he gives rules which (in modern language) are equivalent to the formula \(x^{m+n} = x^{m}x^{n}\) for small values (positive or negative) of m and n ({[}3{]}, vol.~I, pp.~8--13); a little later (pp.~12--13), the ``rule of signs'' is stated, the beginnings of calculating with negative numbers‡; finally, Diophantus uses for the first time a letter to represent an unknown in an equation. In contrast, he seems little concerned to attach a general significance to the methods he applies for the solution of his problems; as for the axiomatic conception of laws of composition as begun by Euclid, this appears foreign to Diophantus' thought as to that of his immediate successors; it only reappears in Algebra at the beginning of the 19th century.

There were first needed, during the intervening centuries, on the one hand the development of a system of algebraic notation adequate for the expression of abstract laws and on the other a broadening of the notion of ``number'' sufficient to give it, through the observation of a wide enough variety of special cases, a general conception. For this purpose, the axiomatic theory of ratios of magnitudes created by the Greeks was insufficient, for it did no more than make precise the intuitive concept of a real number \textgreater0 and the operations on these numbers, which were already known to the Babylonians in a more confused form; now ``numbers'' would be considered of which the Greeks had had no concept and which to begin with had no sensible ``representation'': on the one hand, zero and negative numbers which appeared in the High Middle Ages in Hindu Mathematics; on the other, imaginary numbers, the creation of Italian algebraists of the 16th century.

With the exception of zero, introduced first as a separating sign before being considered as a number (see the Historical Note to Set Theory, III), the common character of these extensions is that (at least to begin with) they were purely ``formal''. By this it must be understood that the new ``numbers'' appeared first of all as the result of operations applied to situations where, keeping to their strict definition, they had no meaning (for example, the difference \(a - b\) of two natural numbers when \(a < b\) ): whence the names ``false'', ``fictive'', ``absurd'', ``impossible'', ``imaginary'', etc., numbers attributed to them. For the Greeks of the Classical Period, enamoured above all of clear thought, such extensions were inconceivable; they could only arise with calculators more disposed than were the Greeks to display a somewhat mystic faith in the power of their methods (``the generality of Analysis'' as the 18th century would say) and to allow themselves to be carried along by the mechanics of their calculations without investigating whether each step was well founded; a confidence moreover usually justified a posteriori, by the exact results to which the extension led, with these new mathematical entities, of rules of calculation uniquely valid, in all rigour, for the numbers which were already known. This explains how little by little these generalizations of the concept of number, which at first occurred only as intermediaries in a sequence of operations whose starting and finishing points were genuine ``numbers'', came to be considered for their own sakes (independent of any application to concrete calculations); once this step had been taken, more or less tangible interpretations were sought, new objects which thus acquired the right to exist in Mathematics.†

On this subject, the Hindus were already aware of the interpretation to be given to negative numbers in certain cases (a debt in a commercial problem, for example). In the following centuries, as the methods and results of Greek and Hindu mathematics spread to the West (via the Arabs), the manipulation of these numbers became more familiar and they took on other ``representations'' of a geometric or kinematic character. With the progressive improvement in algebraic notation, this was the only notable progress in Algebra during the close of the Middle Ages.

At the beginning of the 16th century, Algebra took on a new impulse thanks to the discovery by mathematicians of the Italian school of the solution ``by radicals'' of equations of the 3rd and then of the 4th degree (of which we shall speak in more detail in the Historical Note to Algebra, V); it is at this point that, in spite of their aversion, they felt compelled to introduce into their calculations imaginary numbers; moreover, little by little confidence is born in the calculus of these ``impossible'' numbers, as in that of negative numbers, even though no ``representation'' of them was conceived for more than two centuries.

On the other hand, algebraic notation was decisively perfected by Viète and Descartes; from the latter onwards, algebraic notation is more or less that which we use today.

From the middle of the 17th to the end of the 18th century, it seems that the vast horizons opened by the creation of Infinitesimal Calculus resulted in something of a neglect of Algebra in general and especially of mathematical reflection on laws of composition and on the nature of real and complex numbers.† Thus the composition of forces and the composition of velocities, well known in Mechanics from the end of the 17th century, had no repercussions in Algebra, even though they already contained the germ of vector calculus. In fact it was necessary to await the movement of ideas which, about 1800, led to the geometric representation of complex numbers (see the Historical Note to General Topology, VIII) to see addition of vectors used in Pure Mathematics.‡

It was round about the same time that for the first time in Algebra the notion of law of composition was extended in two different directions to elements which have only distant analogies with ``numbers'' (in the broadest sense of the word up till then). The first of these extensions is due to C. F. Gauss occasioned by his arithmetical researches on quadratic forms \(ax^{2} + bxy + cy^{2}\) with integer coefficients. Lagrange had defined on the set of forms with the same discriminant an equivalence relation† and had on the other hand proved an identity which provided this set with a commutative law of composition (not everywhere defined); starting from these results, Gauss showed that this law is compatible (in the sense of § 1, no. 6) with a refined form of the above equivalence relation ({[}5{]}, vol.~I, p.~272): ``This shows'', he then says, ``what must be understood by the composite class of two or several classes''. He then proceeds to study the ``quotient'' law, establishes essentially that it is (in modern language) a commutative group law and does so with arguments whose generality for the most part goes far beyond the special case considered by Gauss (for example, the argument by which he proves the uniqueness of the inverse element is identical with the one we have given for an arbitrary law of composition in § 2, no 3 (ibid., p.~273)). Nor does he stop there: returning a little later to this question, he recognizes the analogy between composition of classes and multiplication of integers modulo a prime number‡ (ibid., p.~371), but proves also that the group of classes of quadratic forms of given discriminant is not always a cyclic group; the indications he gives on this subject show clearly that he had recognized, at any rate in this particular case, the general structure of finite Abelian groups, which we shall study in Chapter VII ({[}5{]}, vol.~I, p.~374 and vol.~II, p.~266).

The other series of research of which we wish to speak also ended in the concept of a group and its definitive introduction into Mathematics: this is the ``theory of substitutions'', the development of the ideas of Lagrange, Vandermonde and Gauss on the solution of algebraic equations. We do not give here the detailed history of this subject (see the Historical Note to Algebra, V); we must refer to the definition by Ruffini and then Cauchy ({[}6{]}, (2), vol.~I, p.~64) of the ``product'' of two permutations of a finite set,§ and the first notions concerning finite permutation groups: transitivity, primitivity, identity element, permutable elements, etc. But these first results remain, as a whole, quite superficial and Évariste Galois must be considered as the true initiator of the theory: having, in his memorable work {[}8{]}, reduced the study of algebraic equations to that of permutation groups associated with them, he took the study of the latter considerably deeper, both with regard to the general properties of groups (Galois was the first to define the notion of normal subgroup and to recognize its importance) and with regard to the determination of groups with particular properties (where the results he obtained are still today considered to be amongst the most subtle in the theory). The first idea of the ``linear representation of groups'' also goes back to Galois† and this fact shows clearly that he possessed the notion of isomorphism of two group structures, independently of their ``realizations''.

However, if the ingenious methods of Gauss and Galois incontestably led them to a very broad conception of the notion of law of composition, they did not develop their ideas particularly on this point and their works had no immediate effect on the evolution of Abstract Algebra.‡ The clearest progress towards abstraction was made in a third direction: after reflecting on the nature of imaginary numbers (whose geometric representation had brought about, at the beginning of the 19th century, a considerable number of works), algebraists of the English school were the first, between 1830 and 1850, to isolate the abstract notion of law of composition and immediately they broadened the field of Algebra by applying this notion to a host of new mathematical entities: the algebra of Logic with Boole (see the Historical Note to Set Theory, IV), vectors and quaternions with Hamilton, general hypercomplex systems, matrices and non-associative laws with Cayley ({[}10{]}, vol.~I, pp.~127 and 301 and vol.~II, pp.~185 and 475). A parallel evolution was pursued independently on the Continent, notably in Vector Calculus (Möbius, Bellavitis), Linear Algebra and hypercomplex systems (Grassmann), of which we shall speak in more detail in the Historical Note to Algebra, III.§

From all this ferment of original and fertile ideas which revitalized Algebra in the first half of the 19th century the latter emerged thoroughly renewed. Previously methods and results had gravitated round the central problem of the solution of algebraic equations (or Diophantine equations in the Theory of Numbers): ``Algebra'', said Serret in the Introduction to his Course of Higher Algebra {[}12{]}, ``is, properly speaking, the Analysis of equations''. After 1850, if the treatises on Algebra still gave preeminence to the theory of equations for some time to come, new research works were no longer dominated by the concern for immediate applications to the solution of numerical equations but were oriented more and more towards what today is considered to be the essential problem of Algebra, the study of algebraic structures for their own sake.

These works fall quite neatly into three main currents which extend respectively the three movements of ideas which we have analysed above and pursue parallel courses without noticeably influencing one another until the last years of the 19th century.†

The first of these arose out of the work of the German school of the 19th century (Dirichlet, Kummer, Kronecker, Dedekind, Hilbert) in building up the theory of algebraic numbers, issuing out of the work of Gauss who was responsible for the first study of this type, that of the numbers \(a + bi\) ( \(a\) and \(b\) rational). We shall not follow the evolution of this theory here (see the Historical Note to Commutative Algebra, VII): we need only mention, for our purposes, the abstract algebraic notions which arose here. Right from the first successors of Gauss, the idea of a field (of algebraic numbers) is fundamental to all the works on this subject (as also to the works of Abel and Galois on algebraic equations); its field of application was enlarged when Dedekind and Weber {[}13{]} modelled the theory of algebraic functions of one variable on that of algebraic numbers. Dedekind {[}14{]} was also responsible for the notion of an ideal which provides a new example of a law of composition between sets of elements; he and Kronecker were responsible for the more and more important role played by commutative groups and modules in the theory of algebraic fields; we shall return to this in the Historical Notes to Algebra, III, V and VII.

We shall also refer in the Historical Notes to Algebra, III and VIII to the development of Linear Algebra and hypercomplex systems which was pursued without introducing any new algebraic notion during the end of the 19th and the beginning of the 20th century, in England (Sylvester, W. Clifford) and America (B. and C. S. Pierce, Dickson, Wedderburn) following the path traced by Hamilton and Cayley, in Germany (Weierstrass, Dedekind, Frobenius, Molien) and in France (Laguerre, E. Cartan) independently of the Anglo-Saxons and using quite different methods.

As far as group theory is concerned, this developed first of all mainly as the theory of finite permutation groups, following the publication of the works of Galois and their diffusion through the works of Serret {[}12{]} and above all the great ``Traité des Substitutions'' of C. Jordan {[}15{]}. The latter there developed, greatly improving them, the works of his predecessors on the particular properties of permutation groups (transitivity, primitivity, etc.), obtaining results most of which have not been superseded since; he made an equally profound study of a very important particular class of groups, linear groups and their subgroups; moreover, he it was who introduced the fundamental notion of a homomorphism of one group into another and also (a little later) that of quotient group and who proved a part of the theorem known as the ``Jordan-Hölder Theorem''.† Finally Jordan was the first to study infinite groups {[}16{]}, which S. Lie on the one hand and F. Klein and H. Poincaré on the other were to develop considerably in two different directions several years later.

In the meantime, what is essential in a group, that is its law of composition and not the nature of the objects which constitute the group, slowly come to be realized (see for example {[}10{]}, vol.~II, pp.~123 and 131 and {[}14{]}, vol.~III, p.~439). However, even the research works on finite abstract groups had for a long time been conceived as the study of permutation groups and only around 1880 was the independent theory of finite groups beginning to develop consciously. We cannot pursue further the history of this theory which is only touched on very superficially in this Treatise; we shall just mention two of the tools which are still today among the most used in the study of finite groups and which both go back to the 19th century: the Sylow theorems‡ on p-groups which date from 1872 {[}17{]} and the theory of characters created in the last years of the century by Frobenius {[}19{]}. We refer the reader who wishes to go deeply into the theory of finite groups and the many difficult problems it raises, to the monographs of Burnside {[}20{]}, Speiser {[}23{]}, Zassenhaus {[}24{]} and Gorenstein {[}27{]}.

Around 1880 too the systematic study of group presentations was begun; previously only presentations of particular groups had been encountered, for example the alternating group \(\mathfrak{A}_5\) in a work of Hamilton {[}9 bis{]} or the monodromy groups of linear differential equations and Riemann surfaces (Schwarz, Klein, Schläfli). W. Dyck was the first {[}18{]} to define (without giving it a name) the free group generated by a finite number of generators, but he was less interested in it for its own sake than as a ``universal'' object allowing a precise definition of a group given ``by generators and relations''. Developing an idea first set forth by Cayley and visibly influenced on the other hand by the works of Riemann's school mentioned above, Dyck described an interpretation of a group of given presentation, where each generator is represented by a product of two inversions with respect to tangent or intersecting circles (see also Burnside {[}20{]}, Chapter XVIII). A little later, after the development of the theory of automorphic functions by Poincaré and his successors, then the introduction by Poincaré of the tools of Algebraic Topology, the first studies of the fundamental group would be on an equal footing with those of group presentations (Dehn, Tietze), the two theories lending one another mutual support. Moreover it was a topologist, J. Nielsen, who in 1924 introduced the term free group in the first deep study of its properties {[}21{]}, almost immediately afterwards E. Artin (always with reference to questions in topology) introduced the notion of the free product of groups and O. Schreier defined more generally the free product with amalgamated subgroups {[}22{]} (cf.~I, § 7, no. 3, Exercise 20). Here again we cannot go into the history of the later developments in this direction, referring the reader to the work {[}26{]} for more details.

There is no further room to speak of the extraordinary success, since the end of the 19th century, of the idea of group (and that of invariant which is intimately related with it) in Analysis, Geometry, Mechanics and Theoretical Physics. An analogous invasion by this notion and the algebraic notions related to it (groups with operators, rings, ideals, modules) into parts of Algebra which until then seemed quite unrelated is the distinguishing mark of the last period of evolution which we retrace here and which ended with the synthesis of the three branches which we have followed above. This unification is above all the work of the German school of the years 1900--1930: begun by Dedekind and Hilbert in the last years of the 19th century, the work of axiomatization of Algebra was vigorously pursued by E. Steinitz, then, starting in 1920, under the impulse of E. Artin, E. Noether and the algebraists of their school (Hasse, Krull, O. Schreier, van der Waerden). The treatise by van der Waerden {[}25{]}, published in 1930, brought together these works for the first time in a single exposition, opening the way and serving as a guide to many research works in Algebra during more than twenty years.

\BourbakiStarSection{Bibliography}

\begin{enumerate}
\def\labelenumi{\arabic{enumi}.}
\item
  O. NEUGEBAUER, Vorlesungen über Geschichte der antiken Mathematik, Vol. I: Vorgriechische Mathematik, Berlin (Springer), 1934.
\item
  Euclidis Elementa, 5 vols., ed.~J. L. Heiberg, Lipsiae (Teubner), 1883--88.
\end{enumerate}

2 bis. T. L. HEATH, The thirteen books of Euclid's Elements . . ., 3 vols., Cambridge, 1908.

\begin{enumerate}
\def\labelenumi{\arabic{enumi}.}
\setcounter{enumi}{2}
\tightlist
\item
  Diophanti Alexandrini Opera Omnia \ldots, 2 vols., ed.~P. Tannery, Lipsiae (Teubner), 1893--95.
\end{enumerate}

3 bis. Diophante d'Alexandrie, trad. P. Ver Eecke, Bruges (Desclée-de Brouwer), 1926.

\begin{enumerate}
\def\labelenumi{\arabic{enumi}.}
\setcounter{enumi}{3}
\item
  G. W. LEIBNIZ, Mathematische Schriften, ed.~C. I. Gerhardt, vol.~V, Halle (Schmidt), 1858.
\item
  C. F. GAUSS, Werke, vols. I (Göttingen, 1870), II (ibid., 1863) and VIII (ibid., 1900).
\item
  A. L. CAUCHY, Œuvres complètes (2), vol.~I, Paris (Gauthier-Villars), 1905.
\item
  N. H. ABEL, Œuvres, 2 vol., ed.~Sylow and Lie, Christiania, 1881.
\item
  E. GALOIS, Écrits et mémoires mathématiques, ed.~R. Bourgne and J. Y. Azra, Paris (Gauthier-Villars), 1962.
\item
  W. R. HAMILTON, Lectures on Quaternions, Dublin, 1853.
\end{enumerate}

9 bis. W. R. HAMILTON, Memorandum respecting a new system of roots of unity, Phil. Mag. (4), 12 (1856), p.~446.

\begin{enumerate}
\def\labelenumi{\arabic{enumi}.}
\setcounter{enumi}{9}
\item
  A. CAYLEY, Collected mathematical papers, vols. I and II, Cambridge (University Press), 1889.
\item
  H. HANKEL, Vorlesungen über die complexen Zahlen und ihre Functionen, 1st part: Theorie der complexen Zahlensysteme, Leipzig (Voss), 1867.
\item
  J. A. SERRET, Cours d'Algèbre supérieure, 3rd ed., Paris (Gauthier-Villars), 1866.
\item
  R. DEDEKIND and H. WEBER, Theorie der algebraischen Funktionen einer Veränderlichen, Crelle's J., 92 (1882), pp.~181--290.
\item
  R. DEDEKIND, Gesammelte mathematische Werke, 3 vols., Braunschweig (Vieweg), 1932.
\item
  C. JORDAN, Traité des substitutions et des équations algébriques, Paris (Gauthier-Villars), 1870 (reprinted, Paris (A. Blanchard), 1957).
\item
  C. JORDAN, Mémoire sur les groupes des mouvements, Ann. di Mat. (2), 2 (1868), pp.~167--215 and 322--345 (= Œuvres, vol.~IV, pp.~231--302, Paris (Gauthier-Villars), 1964).
\item
  L. SYLOW, Théorèmes sur les groupes de substitutions, Math. Ann., 5 (1872), pp.~584--594.
\item
  W. DYCK, Gruppentheoretische Studien, Math. Ann., 20 (1882), pp.~1-44.
\item
  G. FROBENIUS, Über Gruppencharaktere, Berliner Sitzungsber., 1896, pp.~985--1021 (= Gesammelte Abhandlungen, ed.~J. P. Serre, Berlin-Heidelberg-New York (Springer), vol.~III (1968), pp.~1--37).
\item
  W. BURNSIDE, Theory of groups of finite order, 2nd ed., Cambridge, 1911.
\item
  J. NIELSEN, Die Isomorphismengruppen der freien Gruppen, Math. Ann., 91 (1924), pp.~169--209.
\item
  O. SCHREIER, Die Untergruppen der freien Gruppe, Abh. Hamb., 5 (1927), pp.~161--185.
\end{enumerate}

\begin{enumerate}
\def\labelenumi{\arabic{enumi}.}
\setcounter{enumi}{22}
\item
  A. SPEISER, Theorie der Gruppen von endlicher Ordnung, 3rd ed., Berlin (Springer), 1937.
\item
  H. ZASSENHAUS, Lehrbuch der Gruppentheorie, vol.~I, Leipzig-Berlin (Teubner), 1937.
\item
  B. L. VAN DER WAERDEN, Moderne Algebra, 2nd ed., vol.~I, Berlin (Springer), 1937; vol.~II (ibid.), 1940.
\item
  W. MAGNUS, A. KARRASS and D. SOLITAR, Combinatorial group theory, New York (Interscience), 1966.
\item
  D. GORENSTEIN, Finite groups, New York (Harper and Row), 1968.
\end{enumerate}

\chapter{Linear Algebra}

This chapter is essentially devoted to the study of a particular type of commutative groups with operators (I, § 4, no. 2) called modules. Certain properties stated in §§ 1 and 2 for modules extend to all commutative groups with operators; these will be indicated as they occur. Moreover it will be seen in Chapter III, § 2, no. 6 that the study of a commutative group with operators is always equivalent to that of a module suitably associated with the group with operators in question.

\section{MODULES}
